\section{Shared results VI: Siegel lattices, masses and period integrals}\label{sec:pf:sh-siegel}

\providecommand{\shsiegelGap}{\mathrm{Gap}}
\providecommand{\shsiegelGG}{\gamma}
\providecommand{\shsiegelGam}{\Gamma}
\providecommand{\shsiegelHS}[1]{\lVert #1\rVert_{\mathrm{HS}}}
\providecommand{\shsiegelO}{\mathcal O}
\providecommand{\shsiegelsd}{\mu^{\mathrm{sd}}}

This section proves twenty-eight statements of the formal development \cite{LeanSource2026} which are used, as
black boxes, by several of the preceding sections of Part~II. They concern the Siegel-type objects attached to a
datum: Hida settings and their term systems, the local data at the infinite places, Gram matrices of
Gelfand--Graev characters, masses and volumes of definite hermitian spaces, covolumes of unipotent radicals and local
volumes at the type places, together with a few general facts (adjoints for positive definite hermitian forms, the
Weil representation of unitary dual pairs, and the archimedean theta correspondence on the towers of real data).

The statements are grouped by topic.
\begin{itemize}
\item \Cref{sec:pf:sh-siegel:gap}: numerical consequences of the standing hypotheses: the gap hypothesis, the node gap
condition, and the closedness of the gap condition along classes of families. They are used in
\cref{sec:pf:thmS,sec:pf:thmS-carrier,sec:pf:thmS-leaves,sec:pf:open0,sec:pf:open0a,sec:pf:rows,sec:pf:bundles}.
\item \Cref{sec:pf:sh-siegel:adj}: existence of adjoints and a characterisation of values of adjoint lifts, for
positive definite hermitian forms; used in \cref{sec:pf:bundles,sec:pf:fs}.
\item \Cref{sec:pf:sh-siegel:theta}: the Weil representation of unitary dual pairs (a published result) and the
archimedean theta correspondence on the towers of real data; used in
\cref{sec:pf:bundles,sec:pf:liftinj,sec:pf:open0a}.
\item \Cref{sec:pf:sh-siegel:inf}: the choices of local data at the infinite places and their Wu--Bessel vectors; used
in \cref{sec:pf:thmS-carrier,sec:pf:thmS-leaves,sec:pf:thmS}.
\item \Cref{sec:pf:sh-siegel:orb}: the closed statements on the orbital expansion of the Poincar\'e kernels (level
subgroups, folding, integrability, absolute convergence, the kernel expansion, vanishing of the degenerate terms, the
Gram expansions and row orthogonality); used in \cref{sec:pf:open0,sec:pf:thmS,sec:pf:thmS-leaves}.
\item \Cref{sec:pf:sh-siegel:gg}: the bound for Gram matrices of Gelfand--Graev characters of $\prod_v\GL_k(\F_{q_v})$;
used in \cref{sec:pf:open0,sec:pf:thmS-leaves}.
\item \Cref{sec:pf:sh-siegel:mass}: volumes and masses of definite hermitian spaces; used in
\cref{sec:pf:stageU,sec:pf:top-shape}.
\item \Cref{sec:pf:sh-siegel:covol}: the covolume of the integral points of the unipotent radicals at infinity for the
reference measures, and the local volumes at the type places; used in \cref{sec:pf:multnc,sec:pf:thmS-leaves}.
\item \Cref{sec:pf:sh-siegel:branch}: the length formula on the second branch; used in \cref{sec:pf:top-shape}.
\end{itemize}
Several of the statements are \emph{closed statements}: the formal development phrases a property of all data as a
proposition (for instance $\sfname{FoldOKC}$, ``folding holds for every Hida setting with normalised level subgroup and
Wu decay'') and proves it once; the corresponding targets are these proofs. In such
cases the statement is given in full and the proof assembles the ingredients, which are cited with a one-sentence
statement when they are proved elsewhere. Where a target is a formalised classical theorem, we state it and give the
reference.

Several targets have very long formal proofs: the reference covolume identity (\cref{pf:sh-siegel:covolref}), the volume of
definite spaces (\cref{pf:sh-siegel:volmass}), the kernel expansion (\cref{pf:sh-siegel:kexp}), the vanishing of the degenerate
terms (\cref{pf:sh-siegel:kppoff}), the archimedean theta correspondence (\cref{pf:sh-siegel:arch}) and the second-branch length
formula (\cref{pf:sh-siegel:secondbranch}). For these we give the complete structure of the argument and the arguments at the
top levels, and cite the deepest computations (whose formal proofs are long verifications) by their formal statements.

\subsection{Conventions}\label{sec:pf:sh-siegel:conv}

We use the conventions of \cref{sec:pf:thmS:conv} and \cref{pf:multnc:def-bcf}. A \emph{datum} $D$ has ranks
$k,m,f,t$ and $r=k+m-1$; $m\ge2$ is even, $k\ge m+3$, $t$ is odd, and $r_S=m/2-1$, so that $r=2r_S+k+1$. Its host is a
CM field $E$ over a totally real field $F$, unramified over $F$ at all finite places, with complex conjugation
$x\mapsto\bar x$; the two \emph{type primes} $g_0\ne g_1$ split completely in $E$, and $g_S=g_0g_1$. For every word $u$ of
$D$, $T_u$ is its level prime. The \emph{Harish-Chandra parameters} of $D$ are rational numbers $b_0(\tau),\dots,
b_{m-1}(\tau)$, one family for each embedding type $\tau$. Size data $L_2$ of $D$ carry a scale $L=L_2.L>1$.

A \emph{Hida setting} $X$ for $D$ (as described before \cref{pf:multnc:def-bcf}) consists of the hermitian space $(E^n,J)$, $n=r$, with
its lattice $L_X$ and unitary group $G$, a hyperbolic system $(e_a,f_a)_{a<r_S}$ and for every word $u$ a vector $z_u$
orthogonal to it with $\langle z_u,z_u\rangle=-T_u$, the open compact subgroup $K=K_X\subseteq G(\A_{F,f})$ stabilising
$L_X\otimes\widehat\OO$, a \emph{level subgroup} $K_2\subseteq K$ with a character $\chi_2$ (defined by congruences
modulo the modulus $g$ of $X$ on a reference frame), a Haar measure $\mu_f$ on $G(\A_{F,f})$, Haar measures at the infinite places, and
discrete series data and \emph{Wu functions} $W_u\colon G(F\otimes\R)\to V_u$ (with values in the finite-dimensional
column spaces $V_u$). A \emph{full system} $Y$ over $X$ adds, for every word $u$, a flag $F_u$, a finite set $\mathcal
R_u$ of row lattices, a column system, fundamental sets of the unipotent groups, and a \emph{typed reduction}: finitely
many residue fields $\F_{q_v}$ at the places above $g_S$, with non-trivial additive characters $\psi_v$. For words
$u,u'$ the full system defines the \emph{term system} $\mathcal T(u,u')=Y.t_O(u,u')$ in the sense of
\cref{pf:thmS:terms}; we use the names introduced there ($\mathrm{op}$, $\mathrm{vec}$, rows, the type map $\iota$
with values in the type group $I=\prod_v\GL_k(\F_{q_v})$, $G$, $\Gamma$, $\Gamma_1$, the sets $\mathcal T_{\mathrm{id}}$,
$\mathcal T_{\mathrm{deg}}$, $\mathcal T_{\mathrm{ng}}$, $\mathcal T_{\mathrm{gen}}$). The pair $(X,Y)$ is
\emph{concretely tied} if the level subgroup is normalised, the level subgroups are neat, the column spaces have the
prescribed dimensions, the main term of $X$ is $D_*$, Wu decay and integrability of the phase function hold, the
measures are tied to an adelic bridge and a standard additive character, $D$ is locally tied, and the trace frame
holds (\cref{sec:pf:thmS:conv}).

%==============================================================================
\subsection{Numerical consequences of the standing hypotheses}\label{sec:pf:sh-siegel:gap}

\begin{definition}[The gap hypothesis]\label{pf:sh-siegel:gapdef}
The \emph{gap hypothesis} $\shsiegelGap(D)$ of a datum $D$ (stated before \cref{pf:baselat2:gap}) is the condition
\begin{equation}\label{pf:sh-siegel:eq:gap}
1801\le r\qquad\text{and}\qquad b_{j+1}(\tau)-b_j(\tau)\ \ge\ 16\,r^{19}\quad\text{for all $\tau$ and all $j$ with
$j+1<m$.}
\end{equation}
\formal{Formalism.L2.SG.ConcreteSettingExists.GapHyp}
\end{definition}

Of the standing hypotheses $\mathrm{SH}(D,L_2;C',C_l,C_g,c_b,C_u,c_T,C_T)$ of \cref{pf:open0:def:standing} we use three
components: $C'm^5fL\le k$, $m^4\le f$, and the gap inequality
\begin{equation}\label{pf:sh-siegel:eq:shgap}
C_l\,r^{19}f^5L\ \le\ b_{j}(\tau)-b_{j-1}(\tau)\qquad(1\le j,\ j+1\le m,\ \text{all }\tau).
\end{equation}
(In the formal statement \eqref{pf:sh-siegel:eq:shgap} reads $C_lr^{19}f^5L\le b^{(1)}(\tau,j+1)-b^{(1)}(\tau,j)$, where
$b^{(1)}(\tau,j)=b_{j-1}(\tau)$ for $1\le j\le m$.)

\begin{proposition}[The gap hypothesis from the standing hypotheses]\label{pf:sh-siegel:gaphyp}
Let $D$ be a datum with size data $L_2$, and let $C',C_l,C_g,c_b,C_u,c_T,C_T$ be reals with $C_l\ge16$ and $C'\ge29$. If
the standing hypotheses $\mathrm{SH}(D,L_2;C',C_l,C_g,c_b,C_u,c_T,C_T)$ hold, then $\shsiegelGap(D)$ holds.
\formallabel{Formalism.L2.SG.ThmSOwn.gap_hypothesis_of_hypotheses}
\end{proposition}

\begin{proof}
Write $L=L_2.L$; then $L>1$ (\lref{Formalism.MainNotions.S.one_lt_scL}: the scale of size data exceeds $1$). Since
$m\ge2$ we have $f\ge m^4\ge16$ and $m^5\ge32$, hence
\[
k\ \ge\ C'm^5fL\ \ge\ 29\cdot32\cdot16\cdot1\ =\ 14848\ \ge\ 1801,
\]
and $r=k+m-1\ge k\ge1801$. Now let $\tau$ and $j$ with $j+1<m$. Applying \eqref{pf:sh-siegel:eq:shgap} with index
$j+1$ (so that $1\le j+1$ and $j+2\le m$) gives $C_lr^{19}f^5L\le b_{j+1}(\tau)-b_j(\tau)$. Since $f^5\ge1$, $L\ge1$ and
$C_l\ge16$, $16r^{19}\le C_lr^{19}\le C_lr^{19}f^5L$, which proves \eqref{pf:sh-siegel:eq:gap}.
\end{proof}

The \emph{cubic node gap} and the \emph{node gap condition} of a datum are those of \cref{sec:pf:open0a:points}. With
$N_m=k+m(t+1)$, which equals $\rank(X')+1$ where $\rank(X')=r+mt$ (\lref{Formalism.BasicNotions.factXm}), the
cubic node gap with constant $h=4N_m^3+2t+6$ says that $b_{j+1}(u,\tau)-b_j(u,\tau)\ge h$ for every word $u$, every
$\tau$ and every $j$ with $j+1<m$, where $b_j(u,\tau)$ are the Harish-Chandra parameters of $u$; the node gap condition
says that there is $N\ge\rank(X')+1$ with $b_{j+1}(\tau)-b_j(\tau)\ge t+2N^3+1$ for all $\tau$ and $j+1<m$.

\begin{lemma}[Node gap from the cubic node gap]\label{pf:sh-siegel:nodegap}
For every datum $D$, the cubic node gap with constant $4N_m^3+2t+6$ implies the node gap condition.
\formallabel{GlobalTheory.Core.Rk.GapOfNodes.harishChandraNodesFarApart_of_cubic_node_gap}
\end{lemma}

\begin{proof}
Take $N=N_m=\rank(X')+1$. The parameters of the base word $u_0$ of $D$ are those of $D$ (a defining property of data),
so the hypothesis at $u_0$ gives $b_{j+1}(\tau)-b_j(\tau)\ge4N^3+2t+6\ge t+2N^3+1$, since $N^3\ge0$ and $t\ge0$.
\end{proof}

\begin{definition}[Eventual gap condition, size condition]\label{pf:sh-siegel:gapevdef}
We use the notions of families, classes and closed statements of
\cref{pf:thmS:closed}. For a family $\mathcal F$ the \emph{eventual gap condition} $\shsiegelGap_{\mathrm{ev}}(\mathcal F)$
is: there is $y$ such that $\shsiegelGap(D_{x,i})$ holds for all $x\ge y$ and all $i\in I_x$. The \emph{size condition}
$\sfname{TuSize}(\mathcal F)$ says that eventually the level primes of the members are bounded below as prescribed by the
size data and $\log T_u\le C_{\mathrm{Tu}}\,m\,L$ for every word $u$. The \emph{residual standing
conditions} of a family are $8C_\Gamma\le C_{k,\mathrm{low}}$, $\sfname{TuSize}(\mathcal F)$ and the eventual
arithmetic conditions $\mathrm{ResidEv}(\mathcal F)$ (eventually the root discriminant of the host is at least $2$ and the
arithmetic conditions of the datum hold).
\formal{Formalism.L2.SG.LoopsTies.GapEv}\ \formal{Formalism.MainNotions.TuSize}\
\formal{Formalism.MainNotions.S.ResidualStandingConditions}
\end{definition}

\begin{lemma}[Standing hypotheses along a family]\label{pf:sh-siegel:firstsection}
If a family $\mathcal F$ satisfies the residual standing conditions, then eventually (for $x\ge x_0$ and all $i\in I_x$)
the standing hypotheses $\mathrm{SH}(D_{x,i},L_{x,i};s_{C'}(\mathcal F),s_{C_l}(\mathcal F),C_{\mathrm{YL}}+1,\tfrac14,
3,\tfrac12,C_{\mathrm{Tu}})$ hold.
\formallabel{Formalism.MainNotions.Pk.SD.first_section_hypotheses_at_family_constants}
\end{lemma}

\begin{proof}
The shared lemma \lref{Formalism.MainNotions.S.standing_hypotheses_of_bounds} derives the standing hypotheses at a
member $(D,L)$ of a family, with constants $(C_{k,\mathrm{low}},s_{C_l},C_{\mathrm{YL}}+1,\frac14,3,\frac12,C_{\mathrm{Tu}})$,
from: the parameter choice and size data of the member; the inequalities $s_{C_l}\le C_{\mathrm{Count}}$ and
$s_{C_l}\le C_{\mathrm{gap}}$ and $0<s_{C_l}$ between the constants of the family
(\lref{Formalism.MainNotions.S.sCl_le_CCount}, \lref{Formalism.MainNotions.S.sCl_le_Cgap}, and \lref{Formalism.MainNotions.S.sCl_pos});
three lower bounds for the quantity $X(L)=x^{\eps/100}$ of the size data, namely $X(L)\ge2s_{C_l}$, $X(L)\ge2$ and
$X(L)\ge5/\theta(L)$ with $\theta(L)=\eps/100$; the bound $8C_\Gamma\le C_{k,\mathrm{low}}$; the arithmetic conditions of
the member; and the level-prime window, which follows from $\sfname{TuSize}$ (\lref{Formalism.MainNotions.S.tuWin_of}).
Let $x_0$ be the maximum of the thresholds of $\mathrm{ResidEv}$ and $\sfname{TuSize}$, of $x_1+1$, and of a
threshold $x_3$ beyond which $x^{\eps/100}\ge\max(2s_{C_l},2,500/\eps)$ (which exists since $\eps>0$). For $x\ge x_0$
all inputs hold; the constants of the parameter choice of a member agree with those of the family, so the
conclusion holds with $C_{k,\mathrm{low}}=s_{C'}(\mathcal F)$.
\end{proof}

\begin{proposition}[The gap condition is closed under $\sfname{TuSize}$]\label{pf:sh-siegel:gapclosed}
Let $\mathcal C$ be a class of families such that every family of $\mathcal C$ satisfies $\sfname{TuSize}$. Then
$\sfname{Closed}\bigl(\mathcal C,\ \mathcal F\mapsto(\mathcal F\in\mathcal C\Rightarrow
\shsiegelGap_{\mathrm{ev}}(\mathcal F))\bigr)$.
\formallabel{Formalism.L2.SG.GapClosed.gapEv_closed_of}
\end{proposition}

We first prove the corresponding statement for the class $\mathcal C^T$ of families of $\mathcal C$ satisfying
$\sfname{TuSize}$.

\begin{lemma}\label{pf:sh-siegel:gaprow}
Let $\mathcal C$ be a class such that $\mathrm{ResidEv}$ is closed on $\mathcal C$. Then $\sfname{Closed}(\mathcal
C^T,\mathcal F\mapsto(\mathcal F\in\mathcal C^T\Rightarrow\shsiegelGap_{\mathrm{ev}}(\mathcal F)))$.
\formallabel{Formalism.L2.SG.ThmSOwn.gapRowT}
\end{lemma}

\begin{proof}
Let $\eps_1$ and, for $\eps\le\eps_1$, the staged threshold $K_0=(K_{\mathrm{pre}},K_{\mathrm{mid}},K_{\mathrm{last}})$ be
given by the closedness of $\mathrm{ResidEv}$. Keep $\eps_1$, $K_{\mathrm{pre}}$ and $K_{\mathrm{mid}}$, and replace the
last threshold by
\[
K'_{\mathrm{last}}(l)=\max\Bigl(K_{\mathrm{last}}(l),\ 29,\ 8\sum_i|l_i|\Bigr).
\]
Let $\mathcal F\in\mathcal C^T$ have constants which are large for $(K',C_H)$. They are then large for $(K_0,C_H)$, so
$\mathrm{ResidEv}(\mathcal F)$ holds. Every entry $c$ of the last list of constants satisfies $c\ge29$ and
$c\ge8\sum_i|l_i|$ for $l=c_{\mathrm{out}}\mathbin{+\!\!+}c_{\mathrm{mid}}(C_H)$; since $C_\Gamma$ is an entry of
$c_{\mathrm{mid}}(C_H)$, $8C_\Gamma\le8|C_\Gamma|\le c$. In particular $s_{C'}(\mathcal F)=C_{k,\mathrm{low}}\ge29$ and
$8C_\Gamma\le C_{k,\mathrm{low}}$, and $s_{C_l}(\mathcal F)$, the minimum of six entries of the last list, is $\ge29\ge16$.
Hence the residual standing conditions hold, and \cref{pf:sh-siegel:firstsection} gives $x_0$ such that for $x\ge x_0$ the
standing hypotheses hold at $D_{x,i}$ with $C'=s_{C'}(\mathcal F)\ge29$ and $C_l=s_{C_l}(\mathcal F)\ge16$.
\Cref{pf:sh-siegel:gaphyp} gives $\shsiegelGap(D_{x,i})$ for $x\ge x_0$, that is, $\shsiegelGap_{\mathrm{ev}}(\mathcal
F)$.
\end{proof}

\begin{proof}[Proof of \cref{pf:sh-siegel:gapclosed}]
$\mathrm{ResidEv}$ holds for every family (\lref{Formalism.L2.SG.ResidEvFree.residEv_closed_free}: the residual
eventual conditions are closed on every class, because they hold for every family). Apply \cref{pf:sh-siegel:gaprow}
and take the same $\eps_1$ and thresholds. If $\mathcal F\in\mathcal C$ has large constants, then $\mathcal F\in\mathcal
C^T$ by hypothesis, so $\shsiegelGap_{\mathrm{ev}}(\mathcal F)$.
\end{proof}


%==============================================================================
\subsection{Adjoints for positive definite hermitian forms}\label{sec:pf:sh-siegel:adj}

In this subsection a \emph{sesquilinear form} on a complex vector space $C$ is a map $(x,y)\mapsto\langle x,y\rangle$
which is linear in $x$ and conjugate-linear in $y$; it is \emph{hermitian} if $\langle x,y\rangle=\overline{\langle
y,x\rangle}$ and \emph{positive definite} if $\Re\langle x,x\rangle>0$ for $x\ne0$.

\begin{lemma}[Existence and uniqueness of adjoints]\label{pf:sh-siegel:adjoint}
Let $A$ and $C$ be complex vector spaces with $C$ finite-dimensional, $\langle\cdot,\cdot\rangle_A$ a sesquilinear form on
$A$, $\langle\cdot,\cdot\rangle_C$ a hermitian positive definite sesquilinear form on $C$, $\iota\colon C\to A$ a linear map
and $u\in A$. There is a unique $x\in C$ such that $\langle y,x\rangle_C=\langle\iota(y),u\rangle_A$ for all $y\in C$.
\formallabel{GlobalTheory.hermitian_adjoint_exists_unique}
\end{lemma}

\begin{proof}
Let $C^\vee=\Hom_\C(C,\C)$ and $g\colon C\to C^\vee$, $g(x)(y)=\langle y,x\rangle_C$; $g$ is additive and
$g(cx)=\bar c\,g(x)$, so it is $\R$-linear. It is injective: if $g(x)=0$ then $\langle x,x\rangle_C=0$, so $x=0$ by
positivity. Since $\dim_\R C=2\dim_\C C=2\dim_\C C^\vee=\dim_\R C^\vee<\infty$, the injective $\R$-linear map $g$ is
surjective. The map $y\mapsto\langle\iota(y),u\rangle_A$ belongs to $C^\vee$; its preimage under $g$ is the required $x$,
and uniqueness is the injectivity of $g$.
\end{proof}

\begin{remark}\label{pf:sh-siegel:adjointrem}
The hermitian symmetry of $\langle\cdot,\cdot\rangle_C$ is a hypothesis of the formal statement but is not used in its
proof; no hypothesis on $\langle\cdot,\cdot\rangle_A$ is needed.
\end{remark}

\begin{lemma}[Orthonormal expansions]\label{pf:sh-siegel:onb}
Let $C$ be a finite-dimensional complex vector space with a hermitian positive definite sesquilinear form. There are
$b_1,\dots,b_n\in C$ with $\langle b_i,b_j\rangle=\delta_{ij}$ and $x=\sum_i\langle x,b_i\rangle\,b_i$ for every $x\in C$.
\formallabel{GlobalTheory.ChainOnto.exists_orthonormal_expansion_of_sesq}
\end{lemma}

\begin{proof}
The form $(x,y)\mapsto\langle y,x\rangle$ is an inner product in the usual convention (conjugate-linear in the first
variable); an orthonormal basis $b_1,\dots,b_n$ for it (Gram--Schmidt) has the stated properties, the expansion being the
expansion of $x$ in this basis.
\end{proof}

\begin{proposition}[Values of adjoint lifts]\label{pf:sh-siegel:adjval}
Let $\mathcal X$ be a set, $W$ and $C$ complex vector spaces with $C$ finite-dimensional, and
$\mathrm{ev}\colon C\times W\to\C^{\mathcal X}$ a bilinear map; write $y(w)=\mathrm{ev}(y,w)$. Let $K\subseteq W$ be a
subspace, $\langle\cdot,\cdot\rangle_W$ and $\langle\cdot,\cdot\rangle_C$ hermitian positive definite sesquilinear
forms on $W$ and $C$, and $\mathrm{pet}\colon\C^{\mathcal X}\times\C^{\mathcal X}\to\C$ any map satisfying the
\emph{Schur relation}
\begin{equation}\label{pf:sh-siegel:eq:schur}
\mathrm{pet}\Bigl(y(w),\ \sum_{i=1}^Ny_i(w_i)\Bigr)=\sum_{i=1}^N\langle y,y_i\rangle_C\,\langle w,w_i\rangle_W
\end{equation}
for all $y,y_1,\dots,y_N\in C$ and $w,w_1,\dots,w_N\in W$. Let $x\in C$, $F\in\C^{\mathcal X}$ and $l\colon W\to\C$ be such
that $F$ lies in the span of the functions $y(w)$ with $y\in C$ and $w\in K$, that
$\mathrm{pet}(y(w),F)=\langle y,x\rangle_C\,l(w)$ for all $y\in C$ and $w\in K$, and that $F\ne0$. Then $x\ne0$ and $F=x(w_0)$
for some $w_0\in K$.
\formallabel{GlobalTheory.ChainOnto.value_of_adjoint_of_mem_span}
\end{proposition}

\begin{proof}
Let $b_1,\dots,b_n$ be as in \cref{pf:sh-siegel:onb} for $\langle\cdot,\cdot\rangle_C$, and
$\Phi\colon W^n\to\C^{\mathcal X}$, $\Phi(u)=\sum_jb_j(u_j)$. For $y\in C$ and $w\in K$,
$y(w)=\sum_jb_j(\langle y,b_j\rangle w)=\Phi\bigl((\langle y,b_j\rangle w)_j\bigr)$, and $\langle y,b_j\rangle w\in K$; as
$\Phi$ is linear, the span of these functions is contained in $\Phi(K^n)$. So $F=\sum_jb_j(u_j)$ with all $u_j\in K$.

For every $k$ and $w\in K$, the hypothesis with $y=b_k$ and the Schur relation give
\begin{equation}\label{pf:sh-siegel:eq:key}
\langle w,u_k\rangle_W=\sum_j\langle b_k,b_j\rangle_C\langle w,u_j\rangle_W=\mathrm{pet}(b_k(w),F)=\langle b_k,x\rangle_C\,l(w).
\end{equation}
Since $F\ne0$, some $u_k\ne0$. Taking $w=u_k$ in \eqref{pf:sh-siegel:eq:key}, $\Re(\langle b_k,x\rangle_Cl(u_k))=
\Re\langle u_k,u_k\rangle_W>0$; hence $a:=\langle b_k,x\rangle_C\ne0$, and in particular $x\ne0$. Put
$w_0=\bar a^{-1}u_k\in K$. For $w\in K$, by \eqref{pf:sh-siegel:eq:key},
$\langle w,w_0\rangle_W=a^{-1}\langle w,u_k\rangle_W=l(w)$. For every $j$ the vector $d_j=u_j-\overline{\langle
b_j,x\rangle_C}\,w_0$ lies in $K$ and $\langle w,d_j\rangle_W=\langle b_j,x\rangle_Cl(w)-\langle b_j,x\rangle_Cl(w)=0$ for
all $w\in K$; taking $w=d_j$ gives $d_j=0$. Therefore, using bilinearity, hermitian symmetry and \cref{pf:sh-siegel:onb},
\[
F=\sum_jb_j\bigl(\overline{\langle b_j,x\rangle_C}\,w_0\bigr)=\Bigl(\sum_j\langle x,b_j\rangle_C\,b_j\Bigr)(w_0)=x(w_0).
\qedhere
\]
\end{proof}

%==============================================================================
\subsection{Weil representations and the archimedean theta correspondence}\label{sec:pf:sh-siegel:theta}

\begin{definition}[Weil data]\label{pf:sh-siegel:weildatum}
Let $D$ be a datum, $\psi$ a standard additive character of $\A_F/F$, $S$ and $T$ hermitian
spaces over $E$ attached to $D$, and $\chi_S,\chi_T$ unitary Hecke characters of $\A_E^\times/E^\times$. Write
$\mathbb W=M_{\dim T\times\dim S}(\A_E)$ and $G_S,G_T$ for the unitary groups of $S,T$. A \emph{Weil datum} for
$(\psi,S,T,\chi_S,\chi_T)$ is a homomorphism $\omega\colon G_S(\A)\times G_T(\A)\to\End_\C(\C^{\mathbb W})$ such that:
$\chi_S,\chi_T$ are Hecke characters whose restrictions to $\A_F^\times$ are $\eta^{\dim S}$, $\eta^{\dim T}$ ($\eta$ the
quadratic character of $E/F$); $\omega$ preserves the space of theta functions on $\mathbb W$ and intertwines the
Heisenberg action; $\omega(\gamma,\gamma')f(0)=f(0)$ for rational points $\gamma,\gamma'$ and theta functions $f$; the
functions $p\mapsto\omega(p)f(w)$ are continuous; $\omega(p)$ preserves the $L^2$-norm over a fundamental domain of the
rational lattice $M_{\dim T\times\dim S}(E)$ for every Haar measure; and at every split idempotent, elements supported
there act on Heisenberg-invariant functionals by a positive constant times the splitting factor built from
$\chi_S,\chi_T$.
\formal{Formalism.IndefR.Datum.IsWeilDatumO}
\end{definition}

\begin{littheorem}[Existence of the Weil representation for unitary dual pairs \cite{Weil1964,Kudla1994,HarrisKudlaSweet1996}]
\label{pf:sh-siegel:weil}
For every datum $D$, standard additive character $\psi$, hermitian spaces $S,T$ of $D$ and unitary Hecke characters
$\chi_S,\chi_T$ whose restrictions to $\A_F^\times$ are $\eta^{\dim S}$ and $\eta^{\dim T}$, there is a Weil datum for
$(\psi,S,T,\chi_S,\chi_T)$.
\formallabel{GlobalTheory.weilRepresentationUnitaryPair}
\end{littheorem}

This is the existence of the global Weil representation of the metaplectic group \cite{Weil1964} together with Kudla's
splitting of the metaplectic cover over $G_S(\A)\times G_T(\A)$ attached to a pair of characters with the stated
restrictions \cite{Kudla1994} (see also \cite[\S1]{HarrisKudlaSweet1996}), realised on theta functions on $\mathbb W$.
The development labels it as a published statement, and also proves it: by
\lref{GlobalTheory.WeilExist.exists_isWeilDatumO}, a Weil datum exists as soon as there is a character
$\tilde\eta$ of the kind used to build the splitting at the unramified places (it is first constructed for the
characters $(\tilde\eta^{\dim S},\tilde\eta^{\dim T})$ and then twisted by the Hecke characters
$\tilde\eta^{-\dim S}\chi_S$, $\tilde\eta^{-\dim T}\chi_T$), and such a character exists by
\lref{GlobalTheory.WeilParts.exists_isEtaTilde_unramified}.

\begin{definition}[The archimedean properties of the real data]\label{pf:sh-siegel:archdef}
The \emph{real data} of a datum $D$ are, at each embedding type $\tau$, the discrete series parameters prescribed by the
Harish-Chandra parameters of $D$: the parameters $(p^A_k(\tau),p^A_{k+1}(\tau))$ and $(p^B_k(\tau),p^B_{k+1}(\tau))$ of the
steps of the theta chain ($0\le k<m$), which concern arbitrary hermitian spaces $V,W$ of $D$ carrying discrete series models
with these parameters; the parameters $p^{\mathrm{def}}_i(\tau)$, $p^{\mathrm{hol}}_i(\tau)$ of the pairs $(V^c_i,W^c_i)$ of a
definite source and a holomorphic target ($i<m$); and the archimedean exponents $a^S_k(\tau),a^T_k(\tau)\in\{-e_k(\tau),
-e_k(\tau)-\frac12\}$ of the splitting characters. For a Weil datum $\omega$ for a pair $(V,W)$, discrete series models
$M_V,M_W$ of the unitary groups at $\tau$, and a Fock vector $f$ at $\tau$ (a non-zero theta function which is a vacuum
vector for the Fock model attached to the Cartan involutions of $M_V,M_W$), the Lie algebras $\mathfrak u(V)_\tau$ and
$\mathfrak u(W)_\tau$ act on the polynomial Fock space spanned by the vectors $P\cdot f$ ($P$ a polynomial in
$\dim V\cdot\dim W$ variables). In these terms, the five properties of the real data are:
\begin{enumerate}[label=(\alph*)]
\item\label{pf:sh-siegel:arch:lift} \emph{Lift law} from $M_V$ to $M_W$: there is a non-zero \emph{joint harmonic}
$\Phi$ (harmonic of the same degree for both Lie algebras, a lowest weight vector for $M_V$ and a highest weight vector for
$M_W$), unique up to scalars; for every unitary action of $\mathfrak u(W)_\tau$ on a Hilbert space $U$ and every non-zero
$\mathfrak u(V)_\tau$-invariant map $q$ from (Fock polynomials)$\otimes$($K$-finite vectors of $M_V$) to $U$ intertwining
the $\mathfrak u(W)_\tau$-actions, the $K$-finite vectors of $M_W$ embed into $U$ onto the span of the image of $q$, $q$ does
not vanish on $\Phi$ tensor a unit highest vector, and $q$ is determined by its values on $\Phi$; and the invariant maps
with values in the $K$-finite vectors of $M_W$ form a line.
\item\label{pf:sh-siegel:arch:zeta} \emph{Zeta law}: for a joint harmonic $\Phi$, a unit highest vector $v$ of $M_W$ and a
Haar measure on the real unitary group of $W$, the matrix coefficient of $v$ and the function
$g\mapsto\langle\omega(g)(\Phi f),\Phi f\rangle\overline{\langle\sigma(g)v,v\rangle}$ are integrable, the integral of the
latter is a positive real number, and it equals $\langle\Phi f,\Phi f\rangle/d$ when $V$ is positive definite at $\tau$ and
$d$ is the formal degree of $M_W$ for the chosen measure.
\item\label{pf:sh-siegel:arch:nolift} \emph{No lower lifts}: for every member of the Witt tower strictly below the target
(resp.\ source), every invariant map from the Fock polynomials tensor the $K$-finite vectors of the source (resp.\ target)
model vanishes.
\item\label{pf:sh-siegel:arch:it} \emph{Intertwiner law} $\mathrm{ItLaw}$: along every chain step, for Weil data with
the prescribed exponents and Fock vectors at all $\tau$, the real theta intertwiner of the pure tensor exists and is unique
up to a scalar.
\item\label{pf:sh-siegel:arch:io} \emph{Lowest $K$-type law} $\mathrm{IoLaw}$: for the pairs $(V^c_i,W^c_i)$, the lowest
$K$-type Fock map exists, is unique up to a scalar of absolute value $1$, and takes harmonic values of a fixed degree.
\end{enumerate}
$\mathrm{ChainLower}(D)$ is \ref{pf:sh-siegel:arch:nolift} for the chain steps (target and source), and
$\mathrm{PairLower}(D)$ is \ref{pf:sh-siegel:arch:nolift} for the pairs (target); $\mathrm{ChainTop}(D)$ is
\ref{pf:sh-siegel:arch:lift} to the target and to the source, and \ref{pf:sh-siegel:arch:zeta} on the source, for the
chain steps; $\mathrm{PairTop}(D)$ is \ref{pf:sh-siegel:arch:lift} in both directions and \ref{pf:sh-siegel:arch:zeta} on
the target for the pairs.
\formal{GlobalTheory.IndefiniteRowSideFullHecke.ThetaSourceAction.ArchOnRecord}
\end{definition}
\begin{theorem}[Archimedean theta correspondence on the towers of real data]\label{pf:sh-siegel:arch}
Assume the classification of discrete series of $\Uni(p,q)$ and their covers \textup{(\ref{pf:open0a:P3})} and the
theory of Harish-Chandra modules \textup{(\ref{pf:open0a:P4})}. Then for every datum $D$: $\mathrm{ChainLower}(D)$ and
$\mathrm{PairLower}(D)$ hold, and if the node gap condition holds for $D$, then $\mathrm{ChainTop}(D)$, $\mathrm{PairTop}(D)$,
$\mathrm{ItLaw}(D)$ and $\mathrm{IoLaw}(D)$ hold.
\formallabel{GlobalTheory.T.archimedeanThetaCorrespondence_towers_real_data}
\end{theorem}

These are the local statements of the theta correspondence for the real unitary dual pairs $(\Uni(p,q),\Uni(r,s))$ in
the Fock model: existence and uniqueness of joint harmonics \cite{KashiwaraVergne1978,Howe1989}, the correspondence of
$(\mathfrak g,K)$-modules \cite{Howe1989}, the theta lifts of discrete series in the stable range and their first
occurrence \cite{Li1990,Paul1998}, and the formal-degree identity for the local zeta integrals. The theorem is their
specialisation to the parameters of the towers of real data, deduced in the development from a package of local
statements for general $(p,q,r,s)$.

\begin{proof}
The formal proof uses the following statements of the development, all of which assume \ref{pf:open0a:P3} and
\ref{pf:open0a:P4}.
\begin{enumerate}[label=(\roman*)]
\item\label{pf:sh-siegel:arch:i} A package $X$ of discrete series of all $\Uni(p,q)$ and their covers with the
prescribed parameters, satisfying Blattner's formula, admissibility, and the archimedean Rallis statements for the Fock
model (first occurrence, up--down lines, \dots) (\lref{GlobalTheory.T.exists_provA_discreteSeries_archimedeanRallis}).
\item\label{pf:sh-siegel:arch:ii} The discrete series models of $D$ are cover discrete series
(\lref{GlobalTheory.T.thetaModels_are_coverDiscreteSeries_of_harishChandra}). With
\ref{pf:sh-siegel:arch:i}, a package is \emph{good} for $D$.
\item\label{pf:sh-siegel:arch:iii} The parameters of every chain step and of every pair are \emph{centred block
parameters}: they are signed concatenations of lists $\alpha,\beta,\gamma,\delta$ of rationals and of the block of
length $t$, centred with respect to the exponents, with all entries of absolute value at least $(t+3)/2$ and with the
gap condition when the node gap condition holds
(\lref{GlobalTheory.T.thetaChainStep_centred_block_parameters} and \lref{GlobalTheory.T.definiteSource_pair_centred_block_parameters}).
\item\label{pf:sh-siegel:arch:iv} For a good package, a model of $D$ whose parameter is a small (resp.\ large) centred
block parameter is tied, through its lowest $K$-type, to the corresponding discrete series of the package
(\lref{GlobalTheory.ArchRec.tie_small} and \lref{GlobalTheory.ArchRec.tie_large}); for the lower members of the Witt towers
there is no invariant theta pairing with these discrete series
(\lref{GlobalTheory.T.no_invariant_theta_pairing_of_no_tensor_intertwiner}, which deduces this from the absence of
tensor intertwiners, \lref{GlobalTheory.T.noThetaTensorIntertwiner_lowerTowerMembers}); and the absence of pairings
transfers to the models (\lref{GlobalTheory.T.noThetaMap_models_of_noPairing_covers}).
\item\label{pf:sh-siegel:arch:v} For chain steps the real theta intertwiner exists uniformly and real theta intertwiners
are proportional (\lref{GlobalTheory.T.exists_uniformRealThetaIntertwiner_chainStep} and \lref{GlobalTheory.T.realThetaIntertwiners_proportional_of_harishChandra}).
\item\label{pf:sh-siegel:arch:vi} For the pairs the lowest $K$-type Fock map exists, is unique up to a phase, and is
harmonic (\lref{GlobalTheory.T.lowestKTypeFockMap_exists_of_discreteSeries}, \lref{GlobalTheory.T.lowestKTypeFockMap_unique_of_discreteSeries}, and \lref{GlobalTheory.T.lowestKTypeFockMap_harmonic_of_discreteSeries}).
\end{enumerate}
Fix $D$ and take a good package. \emph{Lower statements}: for a chain step $k$ and $\tau$, write the parameters as in
\ref{pf:sh-siegel:arch:iii}; by \ref{pf:sh-siegel:arch:iv} the source model is tied to the discrete series of the small
parameter, the target models have sign-compatible numberings, and the discrete series admit no invariant pairing with the
Fock spaces of the lower tower members; hence every invariant map from the Fock polynomials of a lower member vanishes,
which is $\mathrm{ChainLower}(D)$ (\lref{GlobalTheory.ArchRec.chainLower}); the same argument with the roles of
small and large exchanged, for the pair parameters, gives $\mathrm{PairLower}(D)$
(\lref{GlobalTheory.ArchRec.pairLower}). \emph{Top statements}: assume the node gap condition. Then the centred
parameters satisfy the gap condition of the Rallis data, and the lift and zeta laws for the chain steps and for the pairs
follow from the Rallis statements of \ref{pf:sh-siegel:arch:i} applied to the tied models
(\lref{GlobalTheory.ArchRec.chainTop} and \lref{GlobalTheory.ArchRec.pairTop}); $\mathrm{ItLaw}(D)$ follows from
\ref{pf:sh-siegel:arch:v} and the chain parameters (\lref{GlobalTheory.ArchRec.itLaw}); and $\mathrm{IoLaw}(D)$ is
the conjunction of the three statements of \ref{pf:sh-siegel:arch:vi} (\lref{GlobalTheory.ArchRec.ioLaw}).
\end{proof}


%==============================================================================
\subsection{Local data at the infinite places}\label{sec:pf:sh-siegel:inf}

Let $X$ be a Hida setting for a datum $D$, and write $N=2r_S+k+1$ ($=n$, \lref{Formalism.L2.SG.Pf.NcReal.JX.nX}).
For an embedding type $\tau$ let $v=v(\tau)$ be the corresponding infinite place of $F$, $J_v$ the diagonal signature
matrix of $X$ at $v$, and $e_a(v),f_a(v),z_u(v)$ the images at $v$ of the hyperbolic system and of the vectors $z_u$;
$w_1(u,v),\dots,w_k(u,v)$ are the orthonormal vectors completing the frame of $u$ at $v$. Let $H_{k,r}$ be the Gram
matrix of the standard frame of size $N$ ($r_S$ hyperbolic pairs, $k$ orthonormal vectors and one vector of norm $-1$;
its signature is $(k+r_S,r_S+1)$) and $\Uni(H_{k,r})$ its unitary group. The discrete series of $X$ at $v$ acts on a Hilbert space $H_v$ with
base vector $w_v$; transported by a matrix $Q$ it gives the representation $\sigma_\tau(m)=\sigma_v(QmQ^{-1})$ of
$\Uni(H_{k,r})$, and similarly $\pi_{u,\tau}$ on the local column space $V_{u,v}$. Put $\lambda_D=\kappa/(2\pi)$, $\kappa$ the
real constant of the host.

\begin{definition}[Choices at infinity]\label{pf:sh-siegel:choice}
Let $\theta\in\R$. A \emph{choice} $\mathcal C$ for $(X,\theta)$ consists of, for every embedding type $\tau$:
\begin{enumerate}[label=(C\arabic*)]
\item\label{pf:sh-siegel:C1} a matrix $Q_\tau\in M_N(\C)$ with inverse $Q_\tau^{-1}$, such that $Q_\tau^*J_vQ_\tau=H_{k,r}$ and the
columns of $Q_\tau$ at the positions of $e_a$ and $f_a$ in the standard basis are $\lambda_D^{r_S-a}e_a(v)$ and
$\overline{\lambda_D^{r_S-a}}^{\,-1}f_a(v)$ ($a<r_S$);
\item\label{pf:sh-siegel:C2} an element $y_\tau\in\Uni(H_{k,r})$ and a Cayley element $P$ such that $\sigma_\tau(y_\tau)w_v$
lies in the lowest $K$-type subspace of the conjugate $g\mapsto\sigma_\tau(P^{-1}gP)$ (for the Harish-Chandra parameter
$\lambda_D(\tau)$ of $D$);
\item\label{pf:sh-siegel:C3} for every word $u$, a Haar probability measure $\mu_W(u,\tau)$ on the stabiliser $\Uni_W(z)$ of the
transported vector $z=z_u(\tau)=Q_\tau^{-1}z_u(v)$;
\item\label{pf:sh-siegel:C4} for every word $u$, a quotient measure $\nu_S(u,\tau)$ of the transported Haar measure
$\mu_G(\tau)$ on $\Uni(H_{k,r})$ by the product $\mu_S(u,\tau)$ of $\mu_W(u,\tau)$ and the Haar measure $\mu_N(\theta)$ of the
unipotent group $N$ normalised by $\theta$.
\end{enumerate}
\formal{Formalism.L2.SG.LocWSOf.Choice}
\end{definition}

\begin{proposition}[Existence of choices]\label{pf:sh-siegel:choiceexists}
Let $X$ be a Hida setting for a datum $D$ with at least one word, and $\theta>0$. Then a choice for $(X,\theta)$ exists.
\formallabel{Formalism.L2.SG.Pf.choiceExists}
\end{proposition}

\begin{lemma}[Transport matrices]\label{pf:sh-siegel:transp}
Let $r,k\ge0$, $N=2r+k+1$, $J$ a diagonal signature matrix of size $N$ with hermitian form $\langle x,y\rangle=x^*Jy$, and
$\lambda\in\C^\times$. Let $e_a,f_a$ ($a<r$), $w_j$ ($j<k$) and $z$ be vectors in $\C^N$ with $\langle e_a,e_b\rangle=
\langle f_a,f_b\rangle=0$, $\langle e_a,f_b\rangle=\delta_{ab}$, $\langle w_i,w_j\rangle=\delta_{ij}$, $\langle z,z\rangle=-1$,
and $e_a,f_a$ orthogonal to $z$ and to the $w_j$, and $z$ orthogonal to the $w_j$. Then there is an invertible $Q\in
M_N(\C)$ with $Q^*JQ=H_{k,r}$ whose columns at the positions of $e_a$ and $f_a$ are $\lambda^{r-a}e_a$ and
$\overline{\lambda^{r-a}}^{\,-1}f_a$. If moreover $\langle z_0,z_0\rangle=-T$ with $T>0$, the vector $z=T^{-1/2}z_0$
satisfies $\langle z,z\rangle=-1$.
\formallabel{Formalism.L2.SG.Pf.Choice.transp_exists}
\formallabel{Formalism.L2.SG.Pf.Choice.transpC_exists}
\end{lemma}

\begin{proof}
The rescaled vectors $\lambda^{r-a}e_a$, $\overline{\lambda^{r-a}}^{\,-1}f_a$, $w_j$, $z$ satisfy the same relations, since
$\langle\lambda^{r-a}e_a,\overline{\lambda^{r-b}}^{\,-1}f_b\rangle=\overline{\lambda^{r-a}}\,\overline{\lambda^{r-b}}^{\,-1}
\delta_{ab}=\delta_{ab}$ (the form is conjugate-linear in the first variable). Let $Q$ be the matrix whose columns are
these vectors, in the order of the standard basis for $H_{k,r}$. Then $Q^*JQ$ is the Gram matrix of the frame, which is
$H_{k,r}$; since $H_{k,r}$ is invertible, so is $Q$, with inverse $H_{k,r}^{-1}Q^*J$. For the last assertion,
$\langle T^{-1/2}z_0,T^{-1/2}z_0\rangle=T^{-1}(-T)=-1$.
\end{proof}

\begin{proof}[Proof of \cref{pf:sh-siegel:choiceexists}]
Fix a word $u_0$. For each $\tau$, the vectors $e_a(v),f_a(v)$, $w_j(u_0,v)$ and $T_{u_0}^{-1/2}z_{u_0}(v)$ satisfy the
hypotheses of \cref{pf:sh-siegel:transp} (the hyperbolic relations of $X$, the frame relations of the vectors $w_j$ and
$\langle z_{u_0},z_{u_0}\rangle=-T_{u_0}$, $T_{u_0}>0$); with $\lambda=\lambda_D\ne0$ this gives \ref{pf:sh-siegel:C1}. For
\ref{pf:sh-siegel:C2}, a Cayley element $P$ for the rows exists (\lref{Formalism.MainNotions.S.Inh.obj_cayleyU_rows}),
and then an element $y_\tau$ of $\Uni(H_{k,r})$ moving the base vector into the lowest $K$-type subspace of the conjugated
representation exists (\lref{Formalism.L2.SG.Pf.Choice.y_exists}, applied to the transported discrete series at $v$,
whose signature and Harish-Chandra parameter are those of $D$ at $\tau$). For \ref{pf:sh-siegel:C3}: the transported vector
$z=z_u(\tau)$ lies in the open cone $W_b$ with $T(z)=T_u>0$ (\lref{Formalism.L2.SG.Pf.Choice.zTr_of}), so its
stabiliser $\Uni_W(z)$ is a compact subgroup (\lref{PI.QuotSetU.UnitaryGroupMassFormula.unitary_weyl_compact}), and a compact
group carries a Haar probability measure (\lref{Formalism.MainNotions.S.Inh.haar_probability_compact_subgroup}). For
\ref{pf:sh-siegel:C4}: $\mu_G(\tau)$ is a Haar measure on $\Uni(H_{k,r})$ (\lref{Formalism.L2.SG.Pf.Choice.haar_tr}),
and the quotient measure exists (\lref{Formalism.L2.SG.Pf.QuotMeasExists.quotMeasExists}). Choosing these data for
all $\tau$ and $u$ gives a choice.
\end{proof}

\begin{definition}[Wu--Bessel data, local factor conditions]\label{pf:sh-siegel:wubdef}
Let $\mathcal C$ be a choice. The \emph{transported Wu function} of $(u,\tau)$ is the function
$w_{u,\tau}$ on $M_N(\C)$ with $w_{u,\tau}(Q_\tau^{-1}\gamma Q_\tau)=W_{u,v}(\gamma)$ for $\gamma$ in the real unitary group at
$v$, and $0$ off $\Uni(H_{k,r})$. \emph{Wu--Bessel data} for $\mathcal C$ are, for every
$(u,\tau)$, a smooth Bessel map $B_{u,\tau}\colon H_v\to V_{u,v}$ for $(z_u(\tau),\psi_S,\sigma_\tau,\pi_{u,\tau})$ such that
$W_{u,v}(\gamma)=B_{u,\tau}(\sigma_v(\gamma)w_v)$ for every $\gamma$. The associated \emph{local Bessel function} is
$\varphi_B(u,\tau)(m)=B_{u,\tau}(\sigma_\tau(m)v_0)$ with $v_0=\sigma_\tau(y_\tau)w_v$. The \emph{local factor conditions}
$\mathrm{LF}(\mathcal C)$ require, for all $(u,\tau)$: $w_v$ is a smooth unit vector; $\mu_S(u,\tau)$ is a Haar measure
on the stabiliser $S(z_u(\tau))=\Uni_W(z_u(\tau))N$; $w_{u,\tau}$ is continuous on $\Uni(H_{k,r})$ and transforms under
$S(z_u(\tau))$ by $\psi_S$ and $\pi_{u,\tau}$; $w_{u,\tau}$ is reproduced by convolution with the normalised matrix
coefficient of $w_v$; and $\int|w_{u,\tau}|^2\,d\nu_S(u,\tau)$ is non-zero and finite.
The \emph{Bessel-of-Whittaker statement} (BW) is the following local assertion, for all $k,r$. Let $\sigma$ be a discrete
series of $\Uni(H_{k,r})$ on a Hilbert space $H$ whose Harish-Chandra parameter has an admissible gap, $w\in H$ a unit vector
in the lowest $K$-type of a Cayley conjugate of $\sigma$ (as in \ref{pf:sh-siegel:C2}), $z\in W_b$ with $T(z)>0$, $c\ne0$,
$\psi_{z,c}$ the associated character of $N$, $\pi$ a finite-dimensional representation of the Harish-Chandra type prescribed
through a frame of $z$, and $\nu$ a quotient of a Haar measure of $\Uni(H_{k,r})$ by
the product of Haar measures of $\Uni_W(z)$ and $N$. Then every function $W$ on $\Uni(H_{k,r})$ which is continuous,
transforms under $S(z)$ by $\psi_{z,c}$ and $\pi$, is reproduced by convolution with the normalised matrix coefficient of
$w$, and has non-zero finite $L^2(\nu)$-norm, is of the form $W(x)=B(\sigma(x)w)$ for a smooth Bessel map $B$ for
$(z,\psi_{z,c},\sigma,\pi)$.
\formal{Formalism.L2.SG.LocWSOf.WuBessel}\ \formal{Formalism.L2.SG.LocWSOf.LocFactor}\
\formal{Formalism.L2.SG.LocWSOf.BesselOfW}
\end{definition}

\begin{proposition}[Wu--Bessel data from the local factors]\label{pf:sh-siegel:wubessel}
Let $\mathcal C$ be a choice for $(X,\theta)$. If $\shsiegelGap(D)$ holds, the Bessel-of-Whittaker statement
(BW) holds, $\mathrm{LF}(\mathcal C)$ holds, every transported representation $\sigma_\tau$ is a discrete
series of $\Uni(H_{k,r})$ with parameter $\lambda_D(\tau)$, every $\mu_G(\tau)$ is a Haar measure and the column types are
transported, then Wu--Bessel data for $\mathcal C$ exist.
\formallabel{Formalism.L2.SG.Pf.wuBesselOfLoc_all}
\end{proposition}

\begin{proof}
This is \lref{Formalism.L2.SG.Pf.WuBL.wuBesselOfLoc}, which derives the conclusion from three statements:
the reproduction of the Wu function under translation (\lref{Formalism.L2.SG.Pf.WuBL.reproduction_translation}),
the invariance of $L^2$-norms for the quotient measure under right translation by $\Uni(H_{k,r})$
(\lref{Formalism.L2.SG.Pf.WuBL.l2RightTranslation_thm}), and the gap bound for the parameters under the gap
hypothesis (\lref{Formalism.L2.SG.Pf.WuBL.parameter_gap_of_hypothesis}). For each $(u,\tau)$, (BW) is applied to
$\sigma_\tau$, the vector $v_0=\sigma_\tau(y_\tau)w_v$ of \ref{pf:sh-siegel:C2}, the transported vector $z_u(\tau)$ and the right
translate $x\mapsto w_{u,\tau}(xy_\tau)$ of the transported Wu function; the three statements turn the conditions
$\mathrm{LF}(\mathcal C)$ into the hypotheses of (BW) for these data (the translate is still reproduced by convolution and has
the same $L^2$-norm, and the parameter has an admissible gap). The resulting Bessel map $B$ satisfies
$w_{u,\tau}(xy_\tau)=B(\sigma_\tau(xy_\tau)w_v)$ on $\Uni(H_{k,r})$, and taking $x=Q_\tau^{-1}\gamma Q_\tau y_\tau^{-1}$
gives $W_{u,v}(\gamma)=B(\sigma_v(\gamma)w_v)$. All three statements are proved in the development, so the implication of
the proposition holds for every choice.
\end{proof}

\begin{proposition}[Local Bessel functions have non-zero finite norm]\label{pf:sh-siegel:l2phi}
Let $\mathcal C$ be a choice satisfying $\mathrm{LF}(\mathcal C)$, with Wu--Bessel data $B$. For every word $u$ and every
$\tau$,
\[
0\ <\ \int\bigl\|\varphi_B(u,\tau)(m)\bigr\|^2\,d\nu_S(u,\tau)(m)\ <\ \infty .
\]
\formallabel{Formalism.L2.SG.Pf.GoodS.l2SG_phiOf_ne_zero_and_lt_infty}
\end{proposition}

\begin{proof}
Write $z=z_u(\tau)$, $w=w_{u,\tau}$, $\nu=\nu_S(u,\tau)$. The vector $z$ lies in $W_b$ with $T(z)=T_u>0$
(\lref{Formalism.L2.SG.Pf.zTr}), $\mu_G(\tau)$ is a Haar measure (\lref{Formalism.L2.SG.Pf.haarTr}), and $\nu$ is a
quotient of $\mu_G(\tau)$ by a Haar measure on $S(z)$; hence $\nu$ is concentrated on $\Uni(H_{k,r})$
(\lref{Formalism.L2.SG.Pf.GoodS.nu_conc}). For $s=h_0n_0\in S(z)$ and $x\in\Uni(H_{k,r})$ the transformation rule in
$\mathrm{LF}(\mathcal C)$ gives $w(sx)=\psi_S(n_0)\,\pi_{u,\tau}(h_0)w(x)$; since $|\psi_S|=1$
(\lref{Formalism.L2.SG.Pf.WuBL.norm_psi}) and $\pi_{u,\tau}(h_0)$ is unitary
(\lref{Formalism.MainNotions.BridgeLemma₃.lbRepOfHC_proved}), $\|w(sx)\|=\|w(x)\|$. The function $w$ is continuous on
$\Uni(H_{k,r})$ and vanishes off it (\lref{Formalism.L2.SG.Pf.WuBL.wT_off}), so the invariance of $L^2$-norms under right
translation (\lref{Formalism.L2.SG.Pf.WuBL.l2RightTranslation_thm}) gives
$\int\|w(xy_\tau)\|^2d\nu=\int\|w(x)\|^2d\nu$. Finally $\varphi_B(u,\tau)(x)=w(xy_\tau)$ for $x\in\Uni(H_{k,r})$
(\lref{Formalism.L2.SG.Pf.GoodS.phi_eq}), hence $\nu$-almost everywhere, so the two integrals agree; the last one is
non-zero and finite by $\mathrm{LF}(\mathcal C)$.
\end{proof}

\begin{remark}
In \cref{pf:sh-siegel:l2phi} the non-vanishing and the finiteness of the norm are hypotheses of the local factor
conditions, for the transported Wu function; the proposition transports them to the local Bessel function.
\end{remark}


%==============================================================================
\subsection{The orbital expansion of the Poincar\'e kernels}\label{sec:pf:sh-siegel:orb}

Throughout this subsection $X$ is a Hida setting for a datum $D$ and $Y$ a full system over $X$. We recall the objects of
the orbital expansion (\cref{pf:thmS:terms} and the description following it).

\begin{definition}[Level subgroups, folding, fundamental sets, orbital integrals]\label{pf:sh-siegel:folddef}
\begin{enumerate}[label=(\alph*)]
\item Let $g$ be the modulus of $X$ and $K(g)=\{x\in K:\ (x-1)L_{X,f}\subseteq g\,L_{X,f}\}$ the principal congruence
subgroup of level $g$ ($L_{X,f}$ the finite adelic completion of the lattice). The level subgroup of $X$ is
\emph{normalised} if every $x\in K(g)$ lies in $K_2$ and satisfies $\chi_2(x)=1$, $K_2\subseteq K$, and
$0<\mu_f(K_2)<\infty$.
\item For a word $u$, $S_{\mathrm{rat}}(u)\subseteq G(\A)$ is the rational subgroup attached to $u$ (the rational points of
the stabiliser of the frame of $u$), and the \emph{phase functions} $\mathrm{Ph}(u,T)$ ($T$ in the column space
$\mathcal C_u$ of $u$) are the functions on $G(\A)$ whose Poincar\'e series over $S_{\mathrm{rat}}(u)\backslash G(F)$ are
the Poincar\'e series of $u$. \emph{Folding holds at $u$}, $\mathrm{Fold}(u)$, if $S_{\mathrm{rat}}(u)$ has a fundamental
domain in $G(\A)$ for $\mu_f$ times the archimedean Haar measure, and there are a compact set $\Omega_0\ni1$ and $a\ge0$ such
that, for every $T\in\mathcal C_u$, the phase function $\mathrm{Ph}(u,T)$ admits fold data with constant $a$ on every compact
$\Omega\supseteq\Omega_0$, and its multiplicity integral over every compact set and every fundamental domain of
$S_{\mathrm{rat}}(u)$ is finite.
\item The terms of the expansion for the pair $(u,u')$ are triples $t=(\Lambda,\Lambda',[\gamma])$ with $\Lambda\in\mathcal
R_u$, $\Lambda'\in\mathcal R_{u'}$ and a class of pairs $(x,\gamma)$ of an adelic point $x$ and a rational isometry $\gamma$
with $\gamma\Lambda'=\Lambda$ which are rationally linked and whose frame $\gamma F_{u'}$ is admissible for $\Lambda$; $p_t$
is the adelic point of a fixed representative. For each term, $R_O(t)\subseteq G(F\otimes\R)$ is the set of archimedean
components of the rationally linked $(x,n)$ with $n$ in the unipotent stabilisers of $(\Lambda,F_u)$ and of
$(\Lambda',F_{u'})$ (after conjugation by $\gamma$). \emph{Fundamental sets exist} for $(u,u')$, $\mathrm{FE}(u,u')$, if every
term $t$ has a strict fundamental set $\Omega$ for $R_O(t)$ in $G(F\otimes\R)$ (measurable, meeting every orbit, and
meeting each orbit at most once); $\Omega_t$ denotes a chosen one.
\item The \emph{orbital integral} of $t$ over a set $\Omega$ is the operator
$\mathrm{OO}_\Omega(t)=\int_\Omega W_u(x)\otimes W_{u'}(\gamma_t^{-1}x)^*\,dx\colon V_{u'}\to V_u$ (an integral of rank-one
operators, $\gamma_t$ the archimedean component of $p_t$); the orbital integrals are \emph{integrable},
$\mathrm{OI}(u,u')$, if the integrand is integrable on $\Omega_t$ for every term $t$. $\mathrm{AC}(u,u')$ is the absolute
convergence of the term system $\mathcal T(u,u')$ (\cref{pf:thmS:terms}).
\end{enumerate}
\formal{Formalism.L2.SG.HidaSetting.HidaSetting.LevelSubgroupNormalized}\
\formal{Formalism.L2.SG.WuDecay.FoldOK}\ \formal{Formalism.L2.SG.TermOp.FundExists}\
\formal{Formalism.L2.SG.TermOp.OOIntAll}\ \formal{Formalism.L2.SG.TermOp.AbsConvT}
\end{definition}

The closed statements are the following: $\sfname{FoldOKC}$: for every $D,X,Y$ and word $u$, if the level subgroup of $X$
is normalised and $X$ has Wu decay, then $\mathrm{Fold}(u)$; $\sfname{FundExistsC}$: for every $D,X,Y,u,u'$,
$\mathrm{FE}(u,u')$; $\sfname{OOIntAllC}$: for every $D,X,Y,u,u'$, $\mathrm{Fold}(u)$, $\mathrm{Fold}(u')$ and
$\mathrm{FE}(u,u')$ imply $\mathrm{OI}(u,u')$; $\sfname{AbsConvC}$: for every $D,X,Y,u,u'$, $\mathrm{Fold}(u)$,
$\mathrm{Fold}(u')$, $\mathrm{FE}(u,u')$ and $\mathrm{OI}(u,u')$ imply $\mathrm{AC}(u,u')$.

\begin{proposition}[Normalised level subgroups]\label{pf:sh-siegel:level}
For every Hida setting $X$ the level subgroup is normalised.
\formallabel{Formalism.L2.SG.Closed.level_subgroup_properties}
\formallabel{Formalism.L2.SG.Pf.levelSubgroupNormalized_of_compact}
\end{proposition}

\begin{proof}
The modulus $g$ is non-zero (\lref{Formalism.L2.SG.Pf.modulus_g_ne_zero}), and the reference vectors
$e_0,f_0,z_0$ and the bases $b_w,b_w^\vee$ entering the definition of $K_2$ and $\chi_2$ lie in $L_{X,f}$ and satisfy the
congruences of the reference frame (\lref{Formalism.L2.SG.Pf.reference_frame_congruences}). Let $x\in K(g)$. Then
$x$ satisfies the congruence conditions defining $K_2$ (\lref{Formalism.L2.SG.Pf.unitaryCongruence.cong_Kg}: $(x-1)$ maps the lattice
into $g$ times the lattice, so it preserves the congruence classes of the reference vectors), so $x\in K_2$. The value
$\chi_2(x)$ is $\psi_{E,f}$ evaluated at $g^{-1}$ times a sum of hermitian products of reference vectors with
$(x-1)$-translates of reference vectors (\lref{Formalism.L2.SG.Pf.unitaryCongruence.flagLevelCharacter_eq_psi_divN}); since
$(x-1)L_{X,f}\subseteq gL_{X,f}$ this argument is integral (\lref{Formalism.L2.SG.Pf.unitaryCongruence.Afun_Kg}), and $\psi_{E,f}$ is
trivial on integral elements (\lref{Formalism.L2.SG.Pf.unitaryCongruence.psi_Z}), so $\chi_2(x)=1$. By definition $K_2\subseteq K$.
Finally, $K(g)$ is open (a property of the adelic instance, as $g\ne0$) and contains $1$, so
$0<\mu_f(K(g))\le\mu_f(K_2)$ (a Haar measure is positive on non-empty open sets), and $K$ is compact
(\lref{PI.WAU.maximal_compact_isCompact}), so $\mu_f(K_2)\le\mu_f(K)<\infty$.
\end{proof}

\begin{proposition}[Folding, fundamental sets, integrability, absolute convergence]\label{pf:sh-siegel:closed4}
The statements $\sfname{FoldOKC}$, $\sfname{FundExistsC}$, $\sfname{OOIntAllC}$ and $\sfname{AbsConvC}$ hold.
\formallabel{Formalism.L2.SG.Closed.FoldOKC_outright}
\formallabel{Formalism.L2.SG.Closed.FundExistsC_holds}
\formallabel{Formalism.L2.SG.Closed.OOIntAllC_outright}
\formallabel{Formalism.L2.SG.Closed.AbsConvC_outright}
\end{proposition}

\begin{proof}
$\sfname{FoldOKC}$ is \lref{Formalism.L2.SG.Closed.FoldOKC_holds} (folding follows from the finiteness of the
Levi quotients, $\sfname{FinLeviC}$, together with the splitting of the level set, the unipotent elements in the level
subgroup, the rational Levi decomposition, the product formula for the character and the unfolding of the integrals)
applied to \lref{Formalism.L2.SG.Closed.FinLeviC_holds}. $\sfname{FundExistsC}$ is
\lref{Formalism.L2.SG.Pf.fundExistsC_holds} (strict fundamental sets exist because the row representatives are
lattices). $\sfname{OOIntAllC}$ and $\sfname{AbsConvC}$ are \lref{Formalism.L2.SG.Closed.OOIntAllC_holds} and
\lref{Formalism.L2.SG.Closed.AbsConvC_holds}, which derive them from the finiteness of the sums of the integrals of the
norms of the orbital integrands over the terms, $\sfname{OOLintSumC}$, which holds by
\lref{Formalism.L2.SG.Closed.OOLintSumC_holds} (a pointwise majorant by columns and a count of slabs).
\end{proof}

\begin{corollary}[Absolute convergence for concretely tied pairs]\label{pf:sh-siegel:absconvtied}
Assume $\sfname{FoldOKC}$, $\sfname{FundExistsC}$, $\sfname{OOIntAllC}$ and $\sfname{AbsConvC}$. Then for every datum $D$,
every concretely tied pair $(X,Y)$ and all words $u,u'$, the term system $\mathcal T(u,u')$ is absolutely convergent.
\formallabel{Formalism.L2.SG.RowsConc.absConvTied_of_full}
\end{corollary}

\begin{proof}
A concrete tie contains the normalisation of the level subgroup and Wu decay, so $\sfname{FoldOKC}$ gives
$\mathrm{Fold}(u)$ and $\mathrm{Fold}(u')$; $\sfname{FundExistsC}$ gives $\mathrm{FE}(u,u')$, then $\sfname{OOIntAllC}$ gives
$\mathrm{OI}(u,u')$, and $\sfname{AbsConvC}$ gives $\mathrm{AC}(u,u')$.
\end{proof}

\subsubsection{The kernel expansion}

\begin{definition}[Kernel forms and term operators]\label{pf:sh-siegel:kexpdef}
Let $u,u'$ be words. For $T\in\mathcal C_u$ and a set $\Delta$ of representatives of $S_{\mathrm{rat}}(u)\backslash G(F)$
the Poincar\'e series is $P_u(T)(x)=\sum_{\delta\in\Delta}\mathrm{Ph}(u,T)(\delta x)$. The \emph{Poincar\'e condition}
$\mathrm{InHH}(u)$ requires that for all $T$ and $\Delta$ this series converges absolutely at every point, does not depend on
$\Delta$, is square integrable on the fundamental domain of $G(F)$ in $G(\A)$, and is fixed by the projection onto the
relevant space. The \emph{kernel form} is
\begin{gather*}
\mathrm{Kf}(T,T')=\bigl(c_0(u)c_0(u')\bigr)^{-1/2}\bigl\langle P_u(T),P_{u'}(T')\bigr\rangle,\\
c_0(u)=\mu_f(K_2)\,\mu_N(u)(\mathrm{Fd}_N(u))\,\mu_H(u)(H_u)/\dim V_u ,
\end{gather*}
with $\mathrm{Fd}_N(u)$ the fundamental set of the unipotent group of $u$ chosen in $Y$. Let $\mathrm{CM}(u)=\bigoplus_\rho V_u$
(the sum over the columns $\rho$ of the column system of $u$, with the $\ell^2$-norm), and $T\mapsto T^\sharp\in\mathrm{CM}(u)$
the map sending $T$ to its values at the column representatives $h_\rho$ (identified with vectors by the Riesz map). For a
term $t$ the \emph{term operator} $a_t\colon\mathrm{CM}(u')\to\mathrm{CM}(u)$ has $(\rho,\rho')$-block
\[
a_t(\rho,\rho')=
\begin{cases}
s_{u,u'}\ \chi_2(k_{\rho\rho'}(t))\ \mathrm{OO}_{\Omega_t}(t)&\text{if $(\Lambda_\rho,\Lambda'_{\rho'})=(\Lambda,\Lambda')$
and $k_{\rho\rho'}(t)\in K_2$,}\\
0&\text{otherwise,}
\end{cases}
\]
with $s_{u,u'}=\bigl(c_0(u)c_0(u')\bigr)^{-1/2}\mu_f(K_2)$, where $k_{\rho\rho'}(t)=(h_\rho x_u)^{-1}p_{t,f}(h_{\rho'}x_{u'})$ is
the finite part of $t$ seen from the columns; the
operator $\mathrm{op}(t)$ of the term system is the image of $a_t$ in the Hilbert--Schmidt space. The \emph{kernel
expansion} $\mathrm{KE}(u,u')$ holds if there is an operator $\mathcal K$ with $\sum_ta_t=\mathcal K$ (unconditional
convergence) and $\mathrm{Kf}(T,T')=\langle\mathcal K(T'^\sharp),T^\sharp\rangle$ for all $T\in\mathcal C_u$, $T'\in\mathcal
C_{u'}$. The \emph{fundamental-set independence} $\mathrm{FI}(u,u')$ requires that $\mathrm{OO}_\Omega(t)=\mathrm{OO}_{\Omega_t}(t)$
for every strict fundamental set $\Omega$ of $R_O(t)$, and the \emph{well-definedness of the amplitudes}
$\mathrm{AW}(u,u')$ requires that the products of $\chi_2$ (and of its factor $\chi_N$) with the orbital integrals do not
depend on the representative of the class of a term.
\formal{Formalism.L2.SG.PoincareInHH.PoincareSetting.InHH}\ \formal{Formalism.L2.SG.PoincareInHH.PoincareSetting.Kform}\
\formal{Formalism.L2.SG.TermOp.aTerm}\ \formal{Formalism.L2.SG.TermOp.KExpansion}
\end{definition}

\begin{theorem}[The kernel expansion]\label{pf:sh-siegel:kexp}
For every datum $D$, Hida setting $X$, full system $Y$ and words $u,u'$: if $\mathrm{Fold}(u)$, $\mathrm{Fold}(u')$,
$\mathrm{InHH}(u)$, $\mathrm{InHH}(u')$, $\mathrm{FE}(u,u')$, $\mathrm{OI}(u,u')$, $\mathrm{FI}(u,u')$, $\mathrm{AW}(u,u')$ and
$\mathrm{AC}(u,u')$ hold, then $\mathrm{KE}(u,u')$ holds.
\formallabel{Formalism.L2.SG.Closed.KExpansionC_holds}
\formallabel{Formalism.L2.SG.Pf.kExpansionC_holds_of}
\formallabel{Formalism.L2.SG.TermOp.KExpCut.kExpansion_of_steps}
\end{theorem}

\begin{proof}
The level subgroup is normalised (\cref{pf:sh-siegel:level}); hence $X$ is \emph{fold-ready}: $G(\A)$ is second countable
and Borel, and the rational points form a discrete subgroup (\lref{Formalism.L2.SG.Pf.RegroupPf.foldReadyX_K2Level}).
Choose:
\begin{enumerate}[label=(\arabic*)]
\item a double transversal $R$ for $S_{\mathrm{rat}}(u)\backslash G(F)/S_{\mathrm{rat}}(u')$
(\lref{Formalism.L2.SG.Pf.Slab.exists_doubleTransversal});
\item for each $r\in R$ a fundamental domain $F_r$ in $G(\A)$ of the stabiliser $S_{\mathrm{rat}}(u)\cap
rS_{\mathrm{rat}}(u')r^{-1}$, which exists because this stabiliser is discrete
(\lref{Formalism.L2.SG.Pf.exists_isFund});
\item a fundamental domain $F_S$ of $S_{\mathrm{rat}}(u)$ (\lref{Formalism.L2.SG.Pf.exists_fund_srat_holds}) and, for
every column $\rho$ of $u$, a strict fundamental set $\Omega_\rho$ of the integral unipotent group of the column
(\lref{Formalism.L2.SG.Pf.Slab.exists_NOcol_fund}).
\end{enumerate}
The phase functions are well defined at $u$ and $u'$ (\lref{Formalism.L2.SG.Pf.phWellDef_holds}). With these choices,
the kernel expansion is obtained in three steps:
\begin{enumerate}[label=(\roman*)]
\item \emph{Entry unfolding} (\lref{Formalism.L2.SG.TermOp.KExpCut.entryUnfold_of}): by folding and the Poincar\'e
conditions, each entry $\mathrm{Kf}(T,T')$ unfolds, through the double cosets $r\in R$ and the fundamental domains $F_r$, into
a sum of integrals of $\mathrm{Ph}(u,T)\,\overline{\mathrm{Ph}(u',T')(r^{-1}\cdot)}$ over $F_r$; the unfolding of a single
double coset is \lref{Formalism.L2.SG.Pf.Slab.double_coset_single_unfolding}.
\item \emph{Double cosets to terms} (\lref{Formalism.L2.SG.TermOp.KExpCut.dcTerms_of}): integrating over the
unipotent fibres (the slabs, using the slab cover property $\sfname{SlabCoverC}$, \lref{Formalism.L2.SG.Closed.SlabCoverC_holds},
neatness of the level subgroups, \lref{Formalism.L2.SG.Closed.NeatKu_holds}, the fundamental bridge,
\lref{Formalism.L2.SG.Pf.fundBridge_holds}, and the measurability of the Poincar\'e series,
\lref{Formalism.L2.SG.Pf.pserMeas_holds}) and regrouping the double cosets according to the lattices and the classes
of their frames (the regrouping statement $\sfname{RegroupC}$, \cref{pf:sh-siegel:regroup} below), the sum of
\textup{(i)} becomes $\sum_t\langle a_t(T'^\sharp),T^\sharp\rangle$.
\item \emph{Summation} (\lref{Formalism.L2.SG.TermOp.KExpCut.kExpansion_of_abs}): by absolute convergence
$\sum_t\|\mathrm{op}(t)\|<\infty$; since the operator norm of a block operator is bounded by a constant times its
Hilbert--Schmidt norm (\lref{Formalism.L2.SG.TermOp.KExpCut.summable_blk_of_abs}), the family $(a_t)$ is summable, with
sum $\mathcal K$, and the identity of (ii) gives $\mathrm{Kf}(T,T')=\langle\mathcal K(T'^\sharp),T^\sharp\rangle$.\qedhere
\end{enumerate}
\end{proof}

\begin{proposition}[Regrouping]\label{pf:sh-siegel:regroup}
The regrouping statement $\sfname{RegroupC}$ holds: for every $D,X,Y,u,u'$ satisfying the hypotheses of
\cref{pf:sh-siegel:kexp} with normalised level subgroup, every family of strict fundamental sets $\Omega_\rho$ of the integral
unipotent groups of the columns $\rho$ of $u$, all $T\in\mathcal C_u$, $T'\in\mathcal C_{u'}$ and every column $\rho$,
\[
s_{u,u'}\int_{\Omega_\rho}\overline{\mathrm{Ph}(u,T)(y,h_\rho x_u)}\,P_{u'}(T')(y,h_\rho x_u)\,dy
=\sum_{\rho'}\Bigl\langle\Bigl(\sum_ta_t(\rho,\rho')\Bigr)(T'^\sharp_{\rho'}),T^\sharp_\rho\Bigr\rangle ,
\]
where the sum over $t$ converges absolutely by $\mathrm{AC}(u,u')$
(\lref{Formalism.L2.SG.TermOp.KExpCut.summable_blk_of_abs}).
\formallabel{Formalism.L2.SG.Pf.RegroupPf.regroupC_all}
\formallabel{Formalism.L2.SG.Pf.RegroupPf.regrouping_of_neatness}
\end{proposition}

\begin{proof}
By \lref{Formalism.L2.SG.Pf.RegroupPf.regrouping_of_closed_lemmas}, regrouping follows from five statements: the phases
are well defined at all words (\lref{Formalism.L2.SG.Pf.RegroupPf.phase_well_defined_all_words}); the column of a
double coset is unique (\lref{Formalism.L2.SG.Pf.RegroupPf.column_uniqueness}); the columns cover the double cosets,
which follows from strong approximation for the nilradical (\lref{Formalism.L2.SG.Pf.RegroupPf.colCoverC_of},
\lref{Formalism.L2.SG.Pf.RegroupPf.nStrongApproxC_of}); absolute convergence over the slabs
(\lref{Formalism.L2.SG.Pf.RegroupPf.slabAbsC_of}, from the slab cover property, fold readiness and neatness); and the
orbit algebra, which identifies the double cosets inside a column with the classes of frames of the terms
(\lref{Formalism.L2.SG.Pf.RegroupPf.orbit_algebra_of_stabilizer}, from the term carrier
\lref{Formalism.L2.SG.Pf.RegroupPf.term_carrier_holds} and the description of stabilisers for neat level
\lref{Formalism.L2.SG.Pf.RegroupPf.stabilizer_of_neat_level}). Fold readiness and neatness hold for normalised level
subgroups (\lref{Formalism.L2.SG.Pf.RegroupPf.foldReadyX_K2Level} and \lref{Formalism.L2.SG.Pf.RegroupPf.neatKu_K2Level}).
\end{proof}

\subsubsection{Degenerate terms, identity terms and row orthogonality}

\begin{definition}[Flags, weights and identity terms]\label{pf:sh-siegel:weightdef}
Let $V$ be a hermitian space over $E$ with form $h$. A \emph{flag} $F$ of length $r$ and norm $T$ is a family
$(e_1,\dots,e_r;z)$ of linearly independent isotropic vectors $e_a$, mutually orthogonal, and a vector $z$ orthogonal to
all $e_a$ with $h(z,z)=-T$; write $F^\natural=(e_1,\dots,e_r,z)$, $X_j(F)=\langle e_1,\dots,e_j\rangle$ and
$X(F)=X_r(F)$. For flags $F,F'$ of norms $T,T'$ let $\widehat Z(F,F')=\bigl(h(F^\natural_a,F'^\natural_b)\bigr)_{a,b}\in
M_{r+1}(E)$, $D_l(F,F')$ its leading principal $l\times l$ minor, and
\[
G(F,F')=D_{r+1}\overline{D_{r+1}}-TT'\,D_r\overline{D_r} .
\]
For the term system $\mathcal T(u,u')$ the value $G(t)=G(F_u,F'_t)$, $F'_t$ the frame of a representative of $t$, does not
depend on the representative (\lref{Formalism.L2.SG.Closed.TermFun_holds}, the first component of which is this
invariance); the weight of $t$ in the term system is $0$ if $G(t)=0$ and $1$ otherwise. For a lattice $\Lambda$, $N_\Lambda(F)$
is the set of automorphisms $n$ of $V$ preserving $\Lambda$ with $(n-1)X_{j+1}(F)\subseteq X_j(F)$ for all $j$ and
$(n-1)X(F)^\perp\subseteq X(F)$. A term $t$ over $\Lambda$ is an \emph{identity term} if some representative frame $F'$
satisfies $F'=nF_u$ modulo $X(F_u)$ for some $n\in N_\Lambda(F_u)$, i.e.\ $n(e_a)=e'_a$ for all $a$ and $n(z)-z'\in X(F_u)$.
\formal{Formalism.MainNotions.S.Gflag}\ \formal{Formalism.L2.SG.FourParts.IsIdF}
\end{definition}

\begin{lemma}\label{pf:sh-siegel:inS}
Let $F,F'$ be flags of length $r$ and norms $T,T'$ with $\bar T=T$ such that $e'_a-e_a\in X_{a-1}(F)$ for $1\le a\le r$ and
$z'-z\in X(F)$. Then $G(F,F')=0$. If $F'=nF$ modulo $X(F)$ for some $n\in N_\Lambda(F)$, then these hypotheses hold.
\formallabel{Formalism.L2.SG.FourParts.gval_eq_zero_of_inSF}
\end{lemma}

\begin{proof}
If $r=0$, then $z'=z$, hence $-T'=h(z',z')=h(z,z)=-T$; so $D_1=h(z,z)=-T$, $D_0=1$ and $G=T\bar T-T^2=0$. If $r\ge1$, then
$e'_1=e_1$ since $X_0(F)=0$. The first column of $\widehat Z(F,F')$ is $(h(F^\natural_a,e_1))_a$, which vanishes: the $e_a$
are isotropic and mutually orthogonal, and $h(z,e_1)=\overline{h(e_1,z)}=0$. Hence $D_l=0$ for all $l\ge1$, in particular
$D_{r+1}=D_r=0$ and $G=0$. The second assertion is the shared lemma \lref{Formalism.L2.SG.FourParts.inSF_of_isIdF}, whose
short proof we include: $e'_a-e_a=(n-1)e_a\in X_{a-1}(F)$, and
$z'-z=(nz-z)-(nz-z')\in X(F)$ because $z\in X(F)^\perp$.
\end{proof}

\begin{proposition}[Identity terms have weight $0$]\label{pf:sh-siegel:iddeg}
For every Hida setting $X$, full system $Y$ and words $u,u'$, the term system $\mathcal T(u,u')$ satisfies $\mathrm{ID}$:
every identity term has weight $0$.
\formallabel{Formalism.L2.SG.Pf.SatGram.idDegenerate_tO}
\formallabel{Formalism.L2.SG.FourParts.G_eq_zero_of_termInS}
\formallabel{Formalism.L2.SG.FourParts.gt_eq_zero_of_isInS}
\end{proposition}

\begin{proof}
Let $t$ be an identity term over $\Lambda\in\mathcal R_u$. The flag $F_u$ has norm $T_u\in\N$, so $\bar T_u=T_u$. By the second
part of \cref{pf:sh-siegel:inS} a representative frame $F'$ of $t$ satisfies the hypotheses of the first part, so
$G(F_u,F')=0$; since $\Lambda$ is a row lattice of $u$ ($Y$ is a full system), $G$ is constant on the class of $t$
(\lref{Formalism.L2.SG.Closed.TermFun_holds}), so $G(t)=0$ and the weight of $t$ is $0$.
\end{proof}

\begin{proposition}[Row orthogonality]\label{pf:sh-siegel:roworth}
For every Hida setting $X$, full system $Y$ and words $u,u'$ (indeed for every choice of the typed reduction, column systems
and fundamental sets), the term system $\mathcal T(u,u')$ satisfies $\mathrm{RO}$.
\formallabel{Formalism.L2.SG.TermOp.rowOrth}
\end{proposition}

\begin{proof}
The row of a term $t=(\Lambda,\Lambda',[\gamma])$ is $(\Lambda,\Lambda',\iota(t))$, and $\mathrm{vec}(t)$ lies in the summand of
$H_1=\bigoplus_{(\Lambda_1,\Lambda_1')}\mathrm{HS}(V_{u'},V_u)$ indexed by $(\Lambda,\Lambda')$: its component at
$(\Lambda_1,\Lambda'_1)$ is $0$ unless $(\Lambda_1,\Lambda'_1)=(\Lambda,\Lambda')$. If $\iota(t)=\iota(t')$ and the rows of $t$
and $t'$ differ, then $(\Lambda,\Lambda')\ne(\Lambda_{t'},\Lambda'_{t'})$, so for every index of the sum defining the inner
product of $H_1$ one of the two components vanishes; hence $\langle\mathrm{vec}(t),\mathrm{vec}(t')\rangle=0$.
\end{proof}

\begin{theorem}[Vanishing of the degenerate terms]\label{pf:sh-siegel:kppoff}
For all reals $C_g,c_b,C_w,c_T,C_T,C_0$ there are $C'_0,C_{l,0}$ (one may take $C'_0=C_{l,0}=0$) such that for all
$C'\ge C'_0$, $C_l\ge C_{l,0}$, every datum $D$ with size data $L_2$, every concretely tied pair $(X,Y)$ and all words
$u\ne u'$: if the $K_2$-hypotheses $\mathrm{H}_K(D,L_2;C',C_l,C_g,c_b,C_w,c_T,C_T,C_0)$ hold, then
$\mathrm{op}(t)=0$ for every degenerate term $t\in\mathcal T_{\mathrm{deg}}(u,u')$.
\formallabel{Formalism.L2.SG.Closed.KppOffC_holds}
\formallabel{Formalism.L2.SG.FixedDecomposition.kppOffC_of_steps0}
\end{theorem}

\begin{proof}
Let $t$ be degenerate; only $G(t)=0$ is used. A concrete tie gives the normalisation of the level subgroup and Wu decay, so
$\mathrm{Fold}(u)$, $\mathrm{Fold}(u')$, $\mathrm{FE}(u,u')$ and $\mathrm{OI}(u,u')$ hold by \cref{pf:sh-siegel:closed4}.
Under $\mathrm{H}_K$, distinct words have distinct parameter vectors $u_w(\tau)\in\R^k$ (the $u$-components of the
Harish-Chandra data of the words) at every embedding type $\tau$
(\lref{Formalism.L2.SG.KappaOffDiagonal.words_differ}). Two cases.

\emph{$t$ is not relevant.} Then $a_t=0$ by \emph{Lemma V} (\lref{Formalism.L2.SG.Closed.LemmaVC_holds}: for terms which are
not relevant, the level character averages to zero over a compact subgroup, so the term operator vanishes).

\emph{$t$ is relevant.} We show that the orbital integral $\mathrm{OO}_{\Omega_t}(t)$ vanishes; then every block of $a_t$
vanishes. By the covariance statement $\sfname{CovRep}_0$ (\lref{Formalism.L2.SG.FixedDecomposition.covariantRepZero_of_lemmaV}),
for every element $h$ of the pointwise stabiliser at infinity of the representative frame of $t$,
\[
\pi_u(h)\circ\mathrm{OO}_{\Omega_t}(t)=\mathrm{OO}_{\Omega_t}(t)\circ\pi_{u'}(\gamma_t^{-1}h\gamma_t).
\]
(Its proof: $h$ and $\gamma_t^{-1}h\gamma_t$ fix the reference vectors of $u$ and $u'$, so they lie in the archimedean
stabilisers, where the archimedean characters are trivial and the Wu functions are equivariant. If the character
$r\mapsto\psi_{\infty,u}(r)\overline{\psi_{\infty,u'}(\gamma_t^{-1}r\gamma_t)}$ is non-trivial on the group containing
$R_O(t)$ cocompactly, the orbital integral vanishes by averaging; otherwise the integrand is invariant under this group, and
the invariance of the measure under shifts by $h$ gives the covariance. The inputs are the averaging lemma, unimodularity at
infinity, the triviality of the level character on the relevant subgroups, and the closedness, discreteness and
cocompactness of the subgroups involved:
\lref{Formalism.L2.SG.Closed.AvgCore_holds}, \lref{Formalism.L2.SG.Pf.ginfUnimod_holds}, \lref{Formalism.L2.SG.Pf.nLamNinf_holds}, \lref{Formalism.L2.SG.Closed.ChiTrivialC_holds}, \lref{Formalism.L2.SG.Pf.psiInfChar_holds}, \lref{Formalism.L2.SG.Pf.rClosedC_holds}, \lref{Formalism.L2.SG.Pf.roDiscreteC_holds}, and \lref{Formalism.L2.SG.TermOp.Vanish.rational_cocompact_holds}.)
By the Schur statement $\sfname{SchurNFC}$ (\lref{schurNFC}: for relevant frames with $G=0$ and words whose
parameters differ at every $\tau$, every operator with this covariance property vanishes, by Schur's lemma for the discrete
series with distinct parameters), $\mathrm{OO}_{\Omega_t}(t)=0$
(\lref{Formalism.L2.SG.FixedDecomposition.OOt_eq_zero0}).

In both cases $\mathrm{op}(t)$, the image of $a_t$ in the Hilbert--Schmidt space, is $0$.
\end{proof}

\subsubsection{The Gram expansions}

\begin{lemma}[Independence of $\chi_N$ of the columns]\label{pf:sh-siegel:chiN}
Let $X$ be a Hida setting and assume the frame tie for $(u,u')$ (the flags $F_u,F_{u'}$ of the full system have the reference
vectors $e_a$ of $X$ and $z_u$, $z_{u'}$). Write $\chi_2=\chi_N\chi_L$ for the factorisation of the level character
(\lref{Formalism.L2.SG.TermOp.level_char_eq_mul}). For every term $t$ and columns $\rho,\rho'$ over the lattices of $t$
with $k_{\rho\rho'}(t)\in K_2$,
\[
\chi_N\bigl(k_{\rho\rho'}(t)\bigr)=\chi_N\bigl(k_{\rho_0\rho_0'}(t)\bigr),
\]
where $\rho_0,\rho'_0$ are the base columns of the lattices of $t$.
\formallabel{Formalism.L2.SG.NormCut.chiNIndep_of_frameTie}
\formallabel{Formalism.L2.SG.NormCut.chiN_colDiff}
\end{lemma}

\begin{proof}
Write $c=(h_\rho x_u)^{-1}(h_{\rho_0}x_u)$ and $c'=(h_{\rho'_0}x_{u'})^{-1}(h_{\rho'}x_{u'})$; then
$k_{\rho\rho'}(t)=c\,k_{\rho_0\rho'_0}(t)\,c'$. The column differences $c,c'$ lie in $K$ and stabilise the reference vectors
modulo $g$ times the lattice (\lref{Formalism.L2.SG.Pf.colLatK}, \lref{Formalism.L2.SG.NormCut.colStab}); the base term
element lies in $K$ and satisfies the congruence conditions of $\chi_N$ (\lref{Formalism.L2.SG.Pf.termK},
\lref{Formalism.L2.SG.NormCut.adm_norm_of_fixed_frames}). On such elements $\chi_N$ is multiplicative
(\lref{Formalism.L2.SG.Pf.ChiNMul.chiN_mul}), so it suffices that $\chi_N(c)=\chi_N(c')=1$. Now $\chi_N(c)$ is
$\psi_{E,f}$ evaluated at $g^{-1}\sum_ah(f_{0,a},c\,\nu_a)$, where $\nu_a$ runs over the reference vectors $e_{0,a+1}$ and
$z_0$; since $c\,\nu_a-\nu_a\in gL_{X,f}$, $f_{0,a}\in L_{X,f}$ and $g^{-1}h(f_{0,a},\nu_a)$ is integral, the argument is
integral, and $\psi_{E,f}$ is trivial on integral elements (\lref{Formalism.L2.SG.Pf.unitaryCongruence.psi_Z}).
\end{proof}

\begin{proposition}[The Gram expansions]\label{pf:sh-siegel:normexp}
For every datum $D$, every concretely tied pair $(X,Y)$ and all words $u,u'$, the term system $\mathcal T(u,u')$ satisfies
$\mathrm{GE}(\mathcal T_{\mathrm{gen}})$ and $\mathrm{GE}(\mathcal T_{\mathrm{ng}})$:
\[
\Bigl\|\sum_{t\in S}\mathrm{op}(t)\Bigr\|^2=\Re\sum_{i,j\in I}\Gamma(i,j)\bigl\langle\mathrm{vp}(S,i),\mathrm{vp}(S,j)\bigr\rangle
\qquad(S=\mathcal T_{\mathrm{gen}},\ \mathcal T_{\mathrm{ng}}).
\]
\formallabel{Formalism.L2.SG.RowsConc.normExpSatTied_holds}
\formallabel{Formalism.L2.SG.NormCut.normExpSatC'_of_PD}
\formallabel{Formalism.L2.SG.NormCut.normExpSatC'_of_acts}
\formallabel{Formalism.L2.SG.NormCut.normExpSatC'_of}
\end{proposition}

\begin{proof}
By \cref{pf:sh-siegel:closed4} and \cref{pf:sh-siegel:absconvtied}, a concrete tie gives $\mathrm{FE}(u,u')$, $\mathrm{Fold}(u)$,
$\mathrm{Fold}(u')$, $\mathrm{OI}(u,u')$ and $\mathrm{AC}(u,u')$; the amplitudes are well defined
(\lref{Formalism.L2.SG.Closed.AmpWellDefC_holds}) and the orbital integrals do not depend on the fundamental set
(\lref{Formalism.L2.SG.Closed.OOFundIndepC_holds}). Under these conditions the Gram expansions follow, by
\lref{Formalism.L2.SG.NormCut.normExpSat_of} (which expands $\|\sum_{t\in S}\mathrm{op}(t)\|^2$ block by block and sums the
phases over the columns), from absolute convergence and the following seven facts.
\begin{enumerate}[label=(\roman*)]
\item The Hilbert--Schmidt inner product of two block operators is the sum over the blocks of the Hilbert--Schmidt inner
products of the blocks (\lref{hilbert_schmidt_blocks}).
\item $\gamma_\psi(x)=|U_I|^{-1}\sum_{y\in I}\prod_v\psi_{N,v}(y_vx_vy_v^{-1})$, where $U_I=\prod_vU_k(\F_v)$ is the unipotent
group of $I$ and $|U_I|=\prod_v|U_k(\F_v)|$ (\lref{Formalism.L2.SG.NormCut.gamma_product_sum_holds}; see
\cref{pf:sh-siegel:ggdef}).
\item For unitriangular $w$: $\psi_N(xw)=\psi_N(x)\psi_N(w)$, $\psi_N(wx)=\psi_N(w)\psi_N(x)$ and $|\psi_N(w)|=1$
(\lref{Formalism.L2.SG.NormCut.psi_natural_equiv_holds}).
\item The $\chi_N$-factor of the blocks is independent of the columns (\cref{pf:sh-siegel:chiN}).
\item \emph{Column phases}: there are labels $\mathrm{sc}(\rho)\in I$ of the columns such that, for columns $\rho,\rho'$ over the
lattices of a term $t$, the value $\chi_L(k_{\rho\rho'}(t))$ (read as $0$ if $k_{\rho\rho'}(t)\notin K_2$) equals
$\psi_I(\mathrm{sc}(\rho)^{-1}\iota(t)\mathrm{sc}(\rho'))$, $\psi_I$ the product of the extended characters $\psi_N$
(\lref{Formalism.L2.SG.NormCut.colPhase_holds}).
\item \emph{Column freeness}: for every row lattice $\Lambda$ there is a subset $A_\Lambda\subseteq I$ such that $(a,\rho,n)\mapsto
a\,\mathrm{sc}(\rho)\,n$ is a bijection from $A_\Lambda\times\{\text{columns over }\Lambda\}\times U_I$ onto $I$, and the row
weight of $\Lambda$ is $|A_\Lambda|^{-1}$ (\lref{Formalism.L2.SG.NormCut.colFree_steps}, from neatness
\lref{Formalism.L2.SG.Closed.NeatC_holds}, the surjectivity of the Levi reduction
\lref{Formalism.L2.SG.NormCut.leviOntoC_holds}, the lattice representatives of the columns
\lref{Formalism.L2.SG.NormCut.colLatRepC} and the masses of the column sets \lref{Formalism.L2.SG.Pf.massAC_holds}).
\item The saturated Gram identity $\sfname{SatGram}$ holds on the generic and on the non-generic terms
(\lref{Formalism.L2.SG.NormCut.satGramC'_of}, from the left and right actions of the type group on the terms, which hold
given the positive definite isomorphism of the data:
\lref{Formalism.L2.SG.NormCut.leftActC'_of_PD}, \lref{Formalism.L2.SG.NormCut.rightActC'_holds}, and \lref{Formalism.L2.SG.Pf.positive_definite_isomorphism}).\qedhere
\end{enumerate}
\end{proof}


%==============================================================================
\subsection{Gram matrices of Gelfand--Graev characters}\label{sec:pf:sh-siegel:gg}

\begin{definition}[Gelfand--Graev characters]\label{pf:sh-siegel:ggdef}
Let $\F$ be a finite field with $q$ elements, $k\in\N$, $G_k=\GL_k(\F)$ and $U_k\subseteq G_k$ the group of upper
unitriangular matrices. For an additive character $\psi_0$ of $\F$ let $\psi_N(u)=\psi_0\bigl(\sum_iu_{i,i+1}\bigr)$ on $U_k$
(a character of $U_k$), extended by $0$ outside $U_k$, and
\[
\gamma_{\psi_0}(x)=\frac1{|U_k|}\sum_{y\in G_k}\psi_N(yxy^{-1})\qquad(x\in G_k),
\]
the character of the Gelfand--Graev representation $\Ind_{U_k}^{G_k}\psi_N$. A \emph{Whittaker functional} on a
representation $(\pi,V)$ of $G_k$ is a linear form $\ell$ with $\ell(\pi(u)v)=\psi_N(u)\ell(v)$ for $u\in U_k$; $\pi$ is
\emph{generic} if it has a non-zero Whittaker functional. For a finite family of finite fields $(\F_v)_{v\in P}$ with
additive characters $\psi=(\psi_v)$ put $I=\prod_v\GL_k(\F_v)$ and $\gamma_\psi(x)=\prod_v\gamma_{\psi_v}(x_v)$ for
$x\in I$.
\formal{Formalism.L2.SG.GGMirabolic.gam}\ \formal{Formalism.L2.SG.GGMirabolic.gamProd}\
\formal{Formalism.L2.SG.GGMirabolic.IsWhit}
\end{definition}

Since $y\mapsto y^{-1}$ permutes $G_k$ and $\psi_N(u^{-1})=\overline{\psi_N(u)}$, we have $\gamma_{\psi_0}(x^{-1})=
\overline{\gamma_{\psi_0}(x)}$, and $\gamma_{\psi_0}(1)=|G_k|/|U_k|>0$ (\lref{Formalism.L2.SG.GGMirabolic.gam_one}).
We use two classical facts, both of which are also proved in the development.

\begin{littheorem}[Multiplicity one of Gelfand--Graev representations \cite{GelfandGraev1962,Zelevinsky1981}]
\label{pf:sh-siegel:ggmult}
If $\psi_0\ne1$ and $\pi$ is an irreducible finite-dimensional representation of $\GL_k(\F)$, then any two Whittaker
functionals on $\pi$ are proportional; that is, $\Ind_{U_k}^{G_k}\psi_N$ is multiplicity free.
\end{littheorem}

In the development this is \lref{Formalism.L2.SG.Closed.GelfandGraevMultOne_holds}, obtained from the multiplicity
one theorem \lref{Formalism.IrreducibleCountingLemmas.GaloisIrreducibility.gelfand_graev_multiplicity_one} (proved by
Gelfand's trick: using the Bruhat decomposition, an anti-involution of $G_k$ compatible with $\psi_N$ is shown to fix every
double coset supporting a $\psi_N$-bi-equivariant function, so that the Hecke algebra of $\psi_N$ is commutative) through
the duality \lref{whittaker_functional_dual}.

\begin{littheorem}[Dimension of generic representations \cite{Gelfand1970,Zelevinsky1981}]\label{pf:sh-siegel:mirab}
If $\psi_0\ne1$ and $\pi$ is an irreducible generic representation of $\GL_k(\F)$, then
$\dim\pi\ge\prod_{i=1}^{k-1}(q^i-1)$.
\end{littheorem}

This follows from the restriction to the mirabolic subgroup: it contains the irreducible representation
$\Ind_{U_k}^{P_k}\psi_N$ of the mirabolic subgroup $P_k$, of dimension $\prod_{i=1}^{k-1}(q^i-1)$. In the development it is
\lref{Formalism.L2.SG.Closed.MirabolicBound_holds} (\lref{mirabolic_bound}).

\begin{lemma}[The convolution scalar]\label{pf:sh-siegel:convscalar}
Let $\psi_0\ne1$ and let $(\pi,V)$ be an irreducible finite-dimensional representation of $G_k=\GL_k(\F)$. Then
\[
\gamma_{\psi_0}(1)^{-1}\sum_{x\in G_k}\gamma_{\psi_0}(x)\,\pi(x)=
\begin{cases}\dfrac{|U_k|}{\dim V}\,\mathrm{id}_V&\text{if $\pi$ is generic,}\\[1ex] 0&\text{otherwise.}\end{cases}
\]
\formallabel{convolution_scalar}
\end{lemma}

\begin{proof}
Substituting $u=yxy^{-1}$,
\begin{align*}
\sum_x\gamma_{\psi_0}(x)\pi(x)&=\frac1{|U_k|}\sum_{y\in G_k}\sum_{u\in U_k}\psi_N(u)\,\pi(y^{-1}uy)=\sum_{y\in G_k}
\pi(y)^{-1}\,\mathcal E\,\pi(y),\\
\mathcal E&=\frac1{|U_k|}\sum_{u\in U_k}\psi_N(u)\pi(u)
\end{align*}
(\lref{PI.GGConvScalarW.representation_conjugate_gamma}). As $\psi_N$ is a character of $U_k$, $\mathcal E$ is an
idempotent (\lref{PI.GGConvScalarW.averaging_operator_idempotent}); by multiplicity one (\cref{pf:sh-siegel:ggmult}) its
rank is at most $1$, and it is $1$ exactly when $\pi$ is generic
(\lref{PI.GGConvScalarW.generic_iff_criterion} and \lref{PI.GGConvScalarW.intertwiner_rank_le_one}). The operator $A=\sum_y
\pi(y)^{-1}\mathcal E\pi(y)$ commutes with $\pi$ (\lref{PI.GGConvScalarW.representation_commutes_gamma}), hence is a
scalar $c$ by Schur's lemma (\lref{PI.GGConvScalarW.schur_lemma_complex}); taking traces,
$c\dim V=|G_k|\tr\mathcal E=|G_k|\rank\mathcal E$ (\lref{PI.GGConvScalarW.trace_representation_gamma} and \lref{PI.GGConvScalarW.idempotent_trace_formula}),
with $\dim V>0$ (\lref{PI.GGConvScalarW.representation_dimension_pos}). Dividing by $\gamma_{\psi_0}(1)=|G_k|/|U_k|$
gives $|U_k|\rank\mathcal E/\dim V$.
\end{proof}

\begin{lemma}[The generic ratio]\label{pf:sh-siegel:genratio}
$|U_k|=q^{k(k-1)/2}$, and for $\psi_0\ne1$ and every irreducible generic representation $\pi$ of $\GL_k(\F)$,
\[
\frac{|U_k|}{\dim\pi}\ \le\ \prod_{i=1}^{k-1}\bigl(1-q^{-i}\bigr)^{-1}.
\]
If $q\ge3$, the right-hand side is at most $1+3/q$.
\formallabel{generic_ratio_bound}
\formallabel{mirabolic_card_unipotent}
\formallabel{Formalism.L2.SG.GGMirabolic.ratio_numeric}
\end{lemma}

\begin{proof}
A unitriangular matrix is determined by its $k(k-1)/2$ entries above the diagonal, which are arbitrary; so
$|U_k|=q^{k(k-1)/2}=\prod_{i=1}^{k-1}q^i$. By \cref{pf:sh-siegel:mirab}, $\dim\pi\ge\prod_{i=1}^{k-1}(q^i-1)>0$, hence
$|U_k|/\dim\pi\le\prod_{i=1}^{k-1}q^i/(q^i-1)=\prod_{i=1}^{k-1}(1-q^{-i})^{-1}$. The last assertion is the numerical
inequality $\prod_{i=1}^{K}(1-q^{-i})^{-1}\le1+3/q$ for $q\ge3$ and all $K$
(\lref{Formalism.FourierDecayAlt.P.GrowthPolynomialOperator.t_GGNumeric}), applied with $K=k-1$.
\end{proof}

\begin{theorem}[The Gram matrix bound]\label{pf:sh-siegel:ggmatrix}
Let $(\F_v)_{v\in P}$ be a finite family of finite fields with $q_v=|\F_v|\ge3$, $\psi_v$ non-trivial additive characters,
$k\in\N$ and $I=\prod_v\GL_k(\F_v)$. Let $\mathcal X$ be a finite set with a free right action of $I$, and
$d\colon\mathcal X\times\mathcal X\to I$ a map with $j\cdot d(i,j)=i$ whenever $i$ and $j$ lie in the same orbit. Let
$M(i,j)=\gamma_\psi(d(i,j))$ if $i,j$ lie in the same orbit and $M(i,j)=0$ otherwise. Then $M$ is positive semi-definite,
and so is
\[
\Bigl(\prod_v\bigl(1+3/q_v\bigr)\Bigr)\,\Re\gamma_\psi(1)\cdot\one-\overline M .
\]
\formallabel{Formalism.L2.SG.Closed.GGMatrixBound_holds}
\formallabel{Formalism.L2.SG.Closed.GGMatrixBoundOfLit_holds}
\formallabel{gg_matrix_bound}
\formallabel{PI.GGMatBoundW.eigBound_of}
\formallabel{PI.GGMatBoundW.local_place_bound}
\end{theorem}

\begin{proof}
Let $\pi$ be the permutation representation of $I$ on $\C^{\mathcal X}$ attached to the action, and
$A=\sum_{x\in I}\gamma_\psi(x)\pi(x)$. By freeness, the only $x$ moving $j$ to $i$ is $d(i,j)$, so the matrix of $A$ in the
standard basis is $M$, and its transpose is $\overline M$ (\lref{PI.GGMatBoundW.matrix_coefficient_equation}, using
$\gamma_\psi(x^{-1})=\overline{\gamma_\psi(x)}$, \lref{PI.GGMatBoundW.gamProd_inv}); in particular $M$ is hermitian
(\lref{PI.GGMatBoundW.totally_hermitian_form}).

\emph{Eigenvalues.} Let $\mu$ be an eigenvalue of $A$. For each $v$ let $A_v=\sum_{x\in\GL_k(\F_v)}\gamma_{\psi_v}(x)\pi(x)$,
where $\GL_k(\F_v)$ is embedded in $I$ as the $v$-component; the $A_v$ commute with each other and with $\pi$, and
$A=\prod_vA_v$. The $\mu$-eigenspace of $A$ contains a non-zero joint eigenvector $x$ of the $A_v$, with eigenvalues $\mu_v$,
and $\mu=\prod_v\mu_v$ (\lref{PI.GGMatBoundW.joint_eigenvector_exists} and \lref{PI.GGMatBoundW.prod_formula'}). Fix $v$.
The restriction of $\pi$ to $\GL_k(\F_v)$ is a sum of irreducible representations; the eigenvector $x$ has a non-zero
component in an irreducible subrepresentation $\sigma$ on which $A_v$ has the eigenvalue $\mu_v$
(\lref{PI.GGMatBoundW.exists_irreducible_le}). By \cref{pf:sh-siegel:convscalar}, $A_v$ acts on $\sigma$ as
$\gamma_{\psi_v}(1)c_v$ with $c_v=|U_k|/\dim\sigma$ if $\sigma$ is generic and $c_v=0$ otherwise, and by
\cref{pf:sh-siegel:genratio}, $0\le c_v\le1+3/q_v$ (\lref{PI.GGMatBoundW.eig_place}). So $\mu_v=\gamma_{\psi_v}(1)c_v$
with $\gamma_{\psi_v}(1)=|\GL_k(\F_v)|/|U_k|>0$, and
\[
\mu=\gamma_\psi(1)\prod_vc_v,\qquad0\le\prod_vc_v\le\prod_v(1+3/q_v),\qquad\gamma_\psi(1)=\Re\gamma_\psi(1)>0 .
\]
\emph{Conclusion.} $A$ is hermitian with eigenvalues in $[0,\prod_v(1+3/q_v)\Re\gamma_\psi(1)]$, so $A$ and
$\prod_v(1+3/q_v)\Re\gamma_\psi(1)-A$ are positive semi-definite (\lref{PI.GGMatBoundW.spectrum_positive_semidefinite});
these are the matrices $M$ and $c\cdot\one-M$, $c=\prod_v(1+3/q_v)\Re\gamma_\psi(1)$, and the transpose $c\cdot\one-\overline M$ of
the latter is positive semi-definite as well.
\end{proof}

\begin{corollary}[The Gram bound for term systems]\label{pf:sh-siegel:termsgg}
Let $(\F_v),\psi,k,I,\mathcal X,d$ be as in \cref{pf:sh-siegel:ggmatrix}. Every term system with type set $\mathcal X$ whose
matrix is $\Gamma(i,j)=\overline{\gamma_\psi(d(i,j))}$ for $i,j$ in the same orbit and $0$ otherwise, and whose constant is
$\Gamma_1=\Re\gamma_\psi(1)$, satisfies $\mathrm{GB}\bigl(\prod_v(1+3/q_v)\bigr)$.
\formallabel{Formalism.L2.SG.Closed.TermsGGBound_holds}
\formallabel{Formalism.L2.SG.GGMirabolic.termsGGBound_of_matrix}
\end{corollary}

\begin{proof}
$\mathrm{GB}(P_g)$ says that $P_g\Gamma_1\cdot\one-\Gamma$ is positive semi-definite; with $P_g=\prod_v(1+3/q_v)$ this is the
second assertion of \cref{pf:sh-siegel:ggmatrix}.
\end{proof}

\begin{proposition}[The Gram bound for the terms of a Hida setting]\label{pf:sh-siegel:ggbound}
Let $X$ be a Hida setting, $Y$ a full system and $u,u'$ words, and let $(\F_v,\psi_v)_{v\in P}$ be the typed reduction of $Y$. If
$\psi_v\ne1$ and $|\F_v|\ge3$ for every $v$, then the term system $\mathcal T(u,u')$ satisfies
$\mathrm{GB}\bigl(\prod_v(1+3/|\F_v|)\bigr)$.
\formallabel{Formalism.L2.SG.TermOp.ggBound}
\end{proposition}

\begin{proof}
The type set of $\mathcal T(u,u')$ is $I=\prod_v\GL_k(\F_v)$, with $\Gamma(i,j)=\overline{\gamma_\psi(j^{-1}i)}$ and
$\Gamma_1=\Re\gamma_\psi(1)$ by definition. Let $I$ act on itself by right multiplication; the action is free and transitive,
and $d(i,j)=j^{-1}i$ satisfies $j\cdot d(i,j)=i$. \Cref{pf:sh-siegel:termsgg} (in the formal statement, the hypothesis
$\sfname{TermsGGBound}$, which holds by \cref{pf:sh-siegel:termsgg}) gives the claim.
\end{proof}


%==============================================================================
\subsection{Volumes and masses of definite hermitian spaces}\label{sec:pf:sh-siegel:mass}

Let $D$ be a datum with host $E/F$. A hermitian space $W$ of $D$ comes with a basis, a distinguished lattice $L_0(W)$ and
signatures $\mathrm{sig}_\tau(W)=(P_\tau,Q_\tau)$ at the embedding types; $W$ is \emph{positive definite} if
$\mathrm{sig}_\tau(W)=(\dim W,0)$ for every $\tau$. The unitary group $U_W$ is viewed over the full ad\`eles $\A_F$, and
$\Gamma_W=U_W(F)$; $K_0(W)\subseteq U_W(\A_F)$ is the stabiliser of the completion $\widehat L_0(W)$ (no condition at the
infinite places).

\begin{definition}[Volumes, masses]\label{pf:sh-siegel:massdef}
\begin{enumerate}[label=(\alph*)]
\item The \emph{mass} of a lattice $L$ in a hermitian space is $\mathrm{mass}(L)=\sum_M|\Aut(M)|^{-1}$, the sum over a set of
representatives $M$ of the isometry classes in the genus of $L$ (in the formal definition: if such a finite set of
representatives exists, and $0$ otherwise). $\mathrm{mass}_n=\mathrm{mass}(\OO_E^n)$ for the standard form; it is positive
(\lref{Formalism.MainNotions.StandardLattice.massI_pos}).
\item $\vol(W)=\vol_D(W)>0$ is the volume attached to $W$ by $D$. A \emph{volume measure} for $W$ of signature $(P,Q)$ is a Haar
measure $\mu$ on $U_W(\A_F)$ such that, for all archimedean embeddings and all frames $A_\rho$ at infinity with
$A_\rho^*J_{P,Q}A_\rho$ equal to the Gram matrix, and all measurable $B_\rho\subseteq\Uni(P,Q)$, the set of $h$ with
$h_f\widehat L_0(W)=\widehat L_0(W)$ and $A_\rho h_\rho A_\rho^{-1}\in B_\rho$ has measure $\prod_\rho c_{P,Q}(B_\rho)$, for fixed
measures $c_{P,Q}$ of $D$ on $\Uni(P,Q)$ (probability measures when $Q=0$). The \emph{volume tie} $\sfname{VolTie}(D)$ says:
for every $W$ of constant signature $(P,Q)$ volume measures exist, and for every volume measure $\mu$ fundamental domains $\mathcal
F$ of $\Gamma_W$ in $U_W(\A_F)$ exist and all of them satisfy $\mu(\mathcal F)=\vol(W)$.
\item $\mathrm{GM}(D,n)$: $\mathrm{mass}(L_0(W))=\mathrm{mass}_n$ for every positive definite $W$ of rank $n$;
$\mathrm{VZ}(D,n)$: $\vol(W)=\mathrm{mass}_n$ for every positive definite $W$ of rank $n$.
\end{enumerate}
\formal{Formalism.HermitianLatticesTheory.mass}\ \formal{Formalism.MainNotions.StandardLattice.VolTie}\
\formal{Formalism.L2.GlobalTheoryTheory.GenusMassVolumes.GenusMassAt}\
\formal{Formalism.L2.GlobalTheoryTheory.GenusMassVolumes.volumeAtLevelZero}
\end{definition}

The two classical inputs are the finiteness of the class number and the identification of the double cosets with the
classes of the genus.

\begin{littheorem}[Finiteness of class numbers \cite{Borel1963}]\label{pf:sh-siegel:classfin}
Let $E/F$ be a CM extension, $J$ a positive definite hermitian matrix over $E$ and $L$ a lattice. Then there are finitely many
$g_1,\dots,g_m\in U_J(\A_F)$ with $U_J(\A_F)=\bigsqcup_i U_J(F)\,g_i\,K_0(L)$.
\end{littheorem}

In the development this is \lref{Formalism.L2.GlobalTheoryTheory.GenusMassVolumes.class_number_finite_rank}
(from \lref{unitary_class_set_finite}, which reduces it, by adelic bounds, to a finiteness statement for classes of
hermitian lattices, \lref{hermitian_lattice_finite}).

\begin{littheorem}[Double cosets and the genus \cite{GanHankeYu2001}]\label{pf:sh-siegel:classmass}
In the situation of \cref{pf:sh-siegel:classfin}, for every such decomposition
\[
\sum_{i=1}^m\frac1{|\Gamma_i|}=\mathrm{mass}(L),\qquad\Gamma_i=\{\gamma\in U_J(F):\ g_i^{-1}\gamma g_i\in K_0(L)\},
\]
because $g\mapsto g\widehat L\cap V$ induces a bijection from $U_J(F)\backslash U_J(\A_F)/K_0(L)$ onto the isometry classes of
the genus of $L$, under which $\Gamma_i$ is the automorphism group of the lattice attached to $g_i$.
\end{littheorem}

In the development this is \lref{Formalism.L2.GlobalTheoryTheory.GenusMassVolumes.classMass_of_D7} (from
\lref{unitary_class_mass_hermitian}, whose proof establishes the bijection: hermitian approximation, lattice
modification at finitely many places, and surjectivity onto the genus).

\begin{proposition}[Volume from class data]\label{pf:sh-siegel:volclass}
Let $D$ be a datum and $n\in\N$. Assume $\sfname{VolTie}(D)$, $\mathrm{GM}(D,n)$, and, for every positive definite $W$ of rank
$n$, the conclusions of \cref{pf:sh-siegel:classfin,pf:sh-siegel:classmass} for $L_0(W)$. Then $\mathrm{VZ}(D,n)$.
\formallabel{Formalism.L2.GlobalTheoryTheory.GenusMassVolumes.volume_mass_of_class_data}
\end{proposition}

\begin{proof}
Let $W$ be positive definite of rank $n$. By the volume tie there are a volume measure $\mu$ and a fundamental domain
$\mathcal F$ of $\Gamma_W$ with $\mu(\mathcal F)=\vol(W)$. As $W$ is definite, the set of the definition of volume measures with
$B_\rho=\Uni(n)$ is $K_0(W)$ (\lref{Formalism.L2.GlobalTheoryTheory.GenusMassVolumes.StabFin.KSub_eq}), so
$\mu(K_0(W))=1$ (\lref{Formalism.L2.GlobalTheoryTheory.GenusMassVolumes.measure_KSub}); $K_0(W)$ is open, hence
measurable. Take the decomposition $U_W(\A_F)=\bigsqcup_i\Gamma_Wg_iK_0(W)$ of \cref{pf:sh-siegel:classfin}. The groups
$\Gamma_i=\Gamma_W\cap g_iK_0(W)g_i^{-1}$ are finite (\cref{pf:sh-siegel:stabfin}), and $\Gamma_W$ is countable
(\lref{Formalism.L2.GlobalTheoryTheory.GenusMassVolumes.Unfold.countableGamma}). Write $\Gamma=\Gamma_W$ and $K=K_0(W)$.
For each $i$, $g_iK$ is stable under left multiplication by $\Gamma_i$, and two points of $g_iK$ lie in the same
$\Gamma$-orbit only if they lie in the same $\Gamma_i$-orbit (if $\gamma g_ik=g_ik'$ then $g_i^{-1}\gamma g_i=k'k^{-1}\in K$). So a
fundamental domain $F_i$ of the finite group $\Gamma_i$ in $g_iK$ is a fundamental domain of $\Gamma$ in the $\Gamma$-stable set
$\Gamma g_iK$, and by left invariance $\mu(K)=\mu(g_iK)=|\Gamma_i|\,\mu(F_i)$. Since fundamental domains of a countable group of
measure preserving maps have equal measure, and the sets $\Gamma g_iK$ partition $U_W(\A_F)$,
\[
\mu(\mathcal F)=\sum_i\mu(\mathcal F\cap\Gamma g_iK)=\sum_i\mu(F_i)=\sum_i\frac{\mu(K)}{|\Gamma_i|}=\sum_i\frac1{|\Gamma_i|}
\]
(\lref{Formalism.L2.GlobalTheoryTheory.GenusMassVolumes.Unfold.measure_fd_toReal}). Hence
$\vol(W)=\sum_i|\Gamma_i|^{-1}=\mathrm{mass}(L_0(W))=\mathrm{mass}_n$, by \cref{pf:sh-siegel:classmass} and $\mathrm{GM}(D,n)$.
\end{proof}

\begin{lemma}[Finite stabilisers]\label{pf:sh-siegel:stabfin}
Let $W$ be a positive definite hermitian space, $K\subseteq U_W(\A_F)$ a subgroup whose elements preserve the completion
$\widehat L_0(W)$, and $g\in U_W(\A_F)$. Then $\{\gamma\in U_W(F):\ g^{-1}\gamma g\in K\}$ is finite.
\formallabel{Formalism.L2.GlobalTheoryTheory.GenusMassVolumes.StabFin.stab_finite}
\end{lemma}

\begin{proof}
Write the elements of $U_W(F)$ as matrices in the basis of $W$. If $g^{-1}\gamma g$ preserves $\widehat L_0(W)$, then a fixed
non-zero multiple $c\in\OO_F$ (depending on $g$ and the basis, not on $\gamma$) of every entry of $\gamma$ is integral
(\lref{Formalism.L2.GlobalTheoryTheory.GenusMassVolumes.StabFin.core}). Since $W$ is positive definite at every embedding,
the entries of elements of $U_W(F)$ are bounded at every complex embedding by a constant depending only on $W$
(\lref{Formalism.L2.GlobalTheoryTheory.GenusMassVolumes.StabFin.unit_bound}). The elements of $c^{-1}\OO_E$ all of whose
conjugates are bounded by a given constant form a finite set; hence there are finitely many such $\gamma$.
\end{proof}

\begin{theorem}[Volume of definite spaces at level zero]\label{pf:sh-siegel:volmass}
Let $D$ be a datum and $n\in\N$. If $\sfname{VolTie}(D)$ and $\mathrm{GM}(D,n)$ hold, then $\vol(W)=\mathrm{mass}_n$ for every
positive definite hermitian space $W$ of $D$ of rank $n$.
\formallabel{Formalism.L2.GlobalTheoryTheory.GenusMassVolumes.volume_mass_of_genus_rank}
\end{theorem}

\begin{proof}
Apply \cref{pf:sh-siegel:volclass}; its class hypotheses hold by \cref{pf:sh-siegel:classfin} and \cref{pf:sh-siegel:classmass}
(\lref{Formalism.L2.GlobalTheoryTheory.GenusMassVolumes.class_number_finite_rank} and \lref{Formalism.L2.GlobalTheoryTheory.GenusMassVolumes.classMass_of_D7}),
applied to the Gram matrix of $W$, which is positive definite at every infinite place, and to $L_0(W)$.
\end{proof}

\begin{definition}[Volume ratios and mass ratios]\label{pf:sh-siegel:vrdef}
The \emph{volume ratio} statement $\mathrm{VR}(D)$ is: for every hermitian space $W$ of $D$
of constant signature $(P,Q)$ with $\dim W\ge3$, and every sign $s\in\{+,-\}$ with $P\ge1$ (if $s=+$) or $Q\ge1$ (if $s=-$),
there is $x\in W$ with $h(x,x)=\pm1$ (sign $s$) such that the orthogonal complement $W'=x^\perp$, with the induced form, is a
hermitian space of $D$ of constant signature $(P-1,Q)$ (resp.\ $(P,Q-1)$) and $\vol(W)/\vol(W')=\mathrm{mr}(\dim W)$, where
$\mathrm{mr}(M)$ is the right-hand side of the mass ratio of \cref{prop:families:mass} (\cref{pf:multnc:eq-mr}). The
\emph{mass ratio} is $\mathrm{mass}_M/\mathrm{mass}_{M-1}=\mathrm{mr}(M)$ for $M\ge3$; given the representation property
$\sfname{HermRepUnram}$ of CM fields (every totally positive element of $F$ is $aN(s)+bN(t)$ for totally positive $a,b$), it
holds by \lref{Formalism.L2.GlobalTheoryTheory.GenusMassVolumes.standardLattice_aux2b} (the Siegel mass ratio).
\formal{Formalism.MainNotions.StandardLattice.Conclusion2}
\end{definition}

\begin{theorem}[The base volume equals the mass]\label{pf:sh-siegel:baseVol}
Assume $\sfname{HermRepUnram}$. Let $D$ be a datum satisfying $\mathrm{VR}(D)$ and $\mathrm{VZ}(D,2)$. Then the base space
$Y_0$ of $D$, of signature $(k+r_S,r_S+1)$ at every $\tau$, satisfies $\vol(Y_0)=\mathrm{mass}_r$.
\formallabel{Formalism.L2.GlobalTheoryTheory.GenusMassVolumes.volume_eq_mass_rank_two}
\formallabel{Formalism.L2.GlobalTheoryTheory.GenusMassVolumes.vol_indef}
\formallabel{Formalism.L2.GlobalTheoryTheory.GenusMassVolumes.volume_mass_ge_two}
\formallabel{Formalism.L2.GlobalTheoryTheory.GenusMassVolumes.volume_mass_succ}
\end{theorem}

\begin{proof}
For a hermitian space of constant signature $(P,Q)$ the rank is $P+Q$ (\lref{Formalism.BasicNotions.sig_sum}).

\emph{Step 1: $\mathrm{VZ}(D,n)$ for all $n\ge2$.} By induction from $n=2$; assume $\mathrm{VZ}(D,n)$ with $n\ge2$ and let $W$ be
positive definite of rank $n+1\ge3$. $\mathrm{VR}(D)$ with $s=+$ gives $W'$ of constant signature $(n,0)$, hence positive
definite of rank $n$, with $\vol(W)/\vol(W')=\mathrm{mr}(n+1)$. By induction $\vol(W')=\mathrm{mass}_n$, and by the mass ratio
$\mathrm{mr}(n+1)=\mathrm{mass}_{n+1}/\mathrm{mass}_n$; since $\mathrm{mass}_n>0$, $\vol(W)=\mathrm{mass}_{n+1}$.

\emph{Step 2: indefinite spaces.} We show by induction on $Q$: every $W$ of constant signature $(P,Q)$ with $P\ge2$ has
$\vol(W)=\mathrm{mass}_{P+Q}$. For $Q=0$, $W$ is positive definite of rank $P\ge2$, and Step 1 applies. For $Q+1$: $\dim
W=P+Q+1\ge3$ and $Q+1\ge1$, so $\mathrm{VR}(D)$ with $s=-$ gives $W'$ of signature $(P,Q)$ with
$\vol(W)/\vol(W')=\mathrm{mr}(P+Q+1)$; by induction $\vol(W')=\mathrm{mass}_{P+Q}$, and the mass ratio gives
$\vol(W)=\mathrm{mass}_{P+Q+1}$.

\emph{Step 3.} The base space $Y_0$ has constant signature $(k+r_S,r_S+1)$ with $k+r_S\ge2$ (as $k\ge5$), so
$\vol(Y_0)=\mathrm{mass}_{k+2r_S+1}=\mathrm{mass}_r$.
\end{proof}

\begin{remark}
\Cref{pf:sh-siegel:baseVol} uses $\mathrm{VZ}(D,n)$ only for $n=2$; the formal statement carries the representation property
$\sfname{HermRepUnram}$ as an instance hypothesis, which the development proves
(\lref{GlobalTheory.hermRepUnram}); it is a consequence of Landherr's theorem \cite{Landherr1936} (a ternary hermitian form
which is indefinite at every real place is isotropic). The formal definition of the mass is $0$ for a lattice whose genus has no
finite set of class representatives; by \cref{pf:sh-siegel:classfin} this case does not occur for definite spaces.
\end{remark}


%==============================================================================
\subsection{Covolumes of unipotent radicals and local volumes}\label{sec:pf:sh-siegel:covol}

\subsubsection{The reference covolume identity}

Let $X$ be a Hida setting for $D$ and write $r=r_S$ in this subsection. For a word $u$ let $N_\infty(u)=\prod_vN_v(u)\subseteq
G(F\otimes\R)$ be the unipotent group attached to the isotropic vectors $e_a$ and to $z_u$ (the group $N_u(F\otimes\R)$ of
\cref{pf:multnc:def-bcf}), and $\Gamma_N(u)\subseteq N_\infty(u)$ its integral points: the archimedean components of the
rational points $\gamma\in G(F)$ with $\gamma_\infty\in N_\infty(u)$, $\gamma_f$ in the finite unipotent group of $u$ and
$x_u^{-1}\gamma_fx_u\in K_2$. A \emph{strict fundamental set} of $\Gamma_N(u)$ is a measurable $\mathrm{Fd}\subseteq
G(F\otimes\R)$ meeting every orbit $\Gamma_N(u)x$ exactly once. A measure $\mu$ on $G(F\otimes\R)$ is \emph{Haar on}
$N_\infty(u)$ if it is $\sigma$-finite, non-zero, concentrated on $N_\infty(u)$ and invariant under left translation by
$N_\infty(u)$. For each infinite place $v$, $\mathrm{ref}_v(u)$ is the \emph{reference measure} on the real unitary group at
$v$ attached to the frame $(e(v),f(v),w(u,v),T_u^{-1/2}z_u(v))$ of $u$; the product
$\mathrm{ref}(u)=\prod_v\mathrm{ref}_v(u)$ is Haar on $N_\infty(w)$ for every word $w$
(\lref{Formalism.L2.SG.Pf.NcReal.JX.refNX_haar}).

\begin{definition}[Coordinates of the unipotent groups]\label{pf:sh-siegel:kapdef}
For $s\in G(F\otimes\R)$, a word $u$ and an infinite place $v$ put, with $\langle\cdot,\cdot\rangle$ the hermitian form at $v$,
\[
x_{ba}=\langle f_b,se_a\rangle,\quad y_b=\langle f_b,sz_u\rangle,\quad\varphi_{bj}=\langle f_b,sw_j(u)\rangle,\quad
A_{ba}=\langle f_b,sf_a\rangle,\quad m=x^{-1}A
\]
($a,b<r$, $j<k$). The \emph{coordinates} $\kappa_u(s)\in\R^{d}$, $d=f(2r^2+2kr+r)$, are, at every $v$: the real and imaginary
parts of $x_{ba}$ ($b<a$), of $y_b$, of $\varphi_{bj}$ and of $\frac12(m_{ba}-\overline{m_{ab}})$ ($b<a$), and $\Im m_{aa}$. The
\emph{unit box} is $\mathrm{Box}(u)=\{s\in N_\infty(u):\ \kappa_u(s)\in[0,1)^d\}$. Put
\begin{gather*}
c(r,k)=-\frac{8r^3+6r^2+6kr^2+r}{6},\\
\mathrm{cov}_N(u)=T_u^{rf}\,2^{-f(rk+r(r-1)/2)}\,|d_F|^{c(r,k)},\qquad\mathrm{two}=2^{fr(r+1)/2}.
\end{gather*}
\formal{Formalism.L2.SG.Pf.LsNc5.kap}\ \formal{Formalism.L2.SG.Pf.LsNc5.boxN}\ \formal{Formalism.L2.SG.Pf.LsNc5.covolN}
\end{definition}

\begin{theorem}[The reference covolume identity]\label{pf:sh-siegel:covolref}
Let $X$ be a Hida setting for a datum $D$ with an adelic bridge, $u,u'$ words and $\mathrm{Fd}_N,\mathrm{Fd}_{N'}$ subsets of
$G(F\otimes\R)$. Assume that the measures of $X$ are tied to the bridge, the level subgroup is normalised, and
$\mathrm{Fd}_N$, $\mathrm{Fd}_{N'}$ are strict fundamental sets of $\Gamma_N(u)$, $\Gamma_N(u')$. Then
\[
\mathrm{ref}(u)(\mathrm{Fd}_N)\cdot\mathrm{ref}(u)(\mathrm{Fd}_{N'})=|d_F|^{e_{DF}},\qquad
e_{DF}=-\frac{8r_S^3+6r_S^2+6kr_S^2+r_S}{3}.
\]
\formallabel{Formalism.L2.SG.Closed.CovolRef_outright}
\formallabel{Formalism.L2.SG.Pf.NcReal.JX.covolRef_of_steps}
\formallabel{Formalism.L2.SG.Pf.NcReal.JX.covolRef_of}
\end{theorem}

The proof has four steps: a covolume in coordinates (\cref{pf:sh-siegel:coordcovol}), its conversion into a covolume for every
Haar measure (\cref{pf:sh-siegel:covolN}), the reference measure of the unit box (\cref{pf:sh-siegel:refbox}), and the final
computation.

\begin{definition}[Block coordinates and block lattices]\label{pf:sh-siegel:blkdef}
Group the coordinates into blocks $p$, with \emph{heights} $h(p)$: the blocks $x_{ba}$
($b<a$, values in $\C^f$, height $a-b$), the row blocks $(\varphi_{b,\cdot},y_b)$ ($b<r$, values in $(\C^{k+1})^f$, height
$2r-b$), the blocks $\frac12(m_{ba}-\overline{m_{ab}})$ ($b<a$, values in $\C^f$, height $2r+1+2r-a-b$), and the blocks
$\Im m_{aa}$ ($a<r$, values in $\R^f$, height $2r+1+2r-2a$). Write $\mathrm{blk}_u(s)=(\mathrm{blk}_u(s)_p)_p$ and $V_p$ for
the space of block $p$, with Lebesgue measure. Let $\mathfrak d=\mathfrak D_F\OO_E$, $\mathfrak D_F$ the different of $F$, and
$\mathfrak a_a=\mathfrak d^{-(r-a)}$. The \emph{block lattices} $\Lambda_p\subseteq V_p$ are generated by the images, under the
archimedean embeddings of $E$, of: $\mathfrak a_b\mathfrak a_a^{-1}$ (blocks $x_{ba}$); the $\OO_E$-module generated by
$\mathfrak a_bL_b$, $L_b$ a lattice of $X$, mapped by $x\mapsto(\langle x,w_j(u)\rangle,\langle x,z_u\rangle)_v$ (row blocks);
$\mathfrak a_b\mathfrak a_a$ (blocks $m_{ba}$); and the imaginary parts of the elements $t\in\mathfrak a_a^2$ with $\bar t=-t$
(blocks $m_{aa}$). Put $\Delta=|d_F|$ and
\begin{gather*}
\mathrm{cv}(x_{ba})=2^{-f}\Delta^{1-2(a-b)},\qquad
\mathrm{cv}(\mathrm{row}_b)=T_u^f\,2^{-f(k+1)}\Delta^{(k+1)(1-2(r-b))},\\
\mathrm{cv}(m_{ba})=2^{-f}\Delta^{1-2(r-b)-2(r-a)},\qquad\mathrm{cv}(m_{aa})=\Delta^{\frac12-2(r-a)} .
\end{gather*}
\formal{Formalism.L2.SG.Pf.NcReal.JX.blk}\ \formal{Formalism.L2.SG.Pf.NcReal.JX.cLat}\ \formal{Formalism.L2.SG.Pf.NcReal.JX.cVal}
\end{definition}

The following four facts are proved in the development by long but direct verifications; we state them as used.
\begin{enumerate}[label=(B\arabic*)]
\item\label{pf:sh-siegel:B1} \emph{Triangular group law} (\lref{Formalism.L2.SG.Pf.NcReal.JX.kapLaw_holds}): there are
measurable maps $B_p$ with $\mathrm{blk}_u(ss')_p=\mathrm{blk}_u(s)_p+\mathrm{blk}_u(s')_p+B_p(\mathrm{blk}_u(s),\mathrm{blk}_u(s'))$
for $s,s'\in N_\infty(u)$, where $B_p(x,y)$ depends only on the blocks of height $<h(p)$ and $B_p(x,0)=0$.
\item\label{pf:sh-siegel:B2} \emph{Sections} (\lref{Formalism.L2.SG.Pf.NcReal.JX.kapSurj_holds}): there is a measurable
$\sigma\colon\prod_pV_p\to N_\infty(u)$ with $\mathrm{blk}_u(\sigma(t))=t$; and $\kappa_u$ is injective on $N_\infty(u)$
(\lref{Formalism.L2.SG.Pf.LsNc5.kapInj}).
\item\label{pf:sh-siegel:B3} \emph{Fibration of the integral points} (\lref{Formalism.L2.SG.Pf.NcReal.JX.NOf.nofib_holds}):
$\Gamma_N(u)$ is a subgroup of $N_\infty(u)$, and there are maps $\tau_p$, depending only on the blocks of height $<h(p)$, such
that $x\in\mathrm{blk}_u(\Gamma_N(u))$ if and only if $x_p-\tau_p(x)\in\Lambda_p$ for all $p$.
\item\label{pf:sh-siegel:B4} \emph{Block covolumes} (\lref{Formalism.L2.SG.Pf.NcReal.JX.CovB.blockCovol_holds}): each
$\Lambda_p$ is countable, has a measurable fundamental domain, and every measurable fundamental domain of $\Lambda_p$ has
Lebesgue measure $\mathrm{cv}(p)$.
\end{enumerate}

\begin{lemma}[Product of the block covolumes]\label{pf:sh-siegel:cvprod}
$\prod_p\mathrm{cv}(p)=\mathrm{cov}_N(u)/\mathrm{two}$.
\formallabel{Formalism.L2.SG.Pf.NcReal.JX.cVal_prod}
\end{lemma}

\begin{proof}
There are $\binom r2$ blocks of each of the types $x_{ba}$, $m_{ba}$ and $r$ of each of the types $\mathrm{row}_b$, $m_{aa}$. The
power of $T_u$ is $T_u^{rf}$. The power of $2$ is $2^{-f(r(r-1)+r(k+1))}=2^{-f(rk+r^2)}$, and $2^{-f(rk+r(r-1)/2)}/\mathrm{two}=
2^{-f(rk+r^2)}$. For the power of $\Delta$, with $j=r-b$ and $\sum_{b<a}(a-b)=\frac{r(r^2-1)}6$,
$\sum_{b<a}\bigl((r-b)+(r-a)\bigr)=(r-1)\frac{r(r+1)}2$, the exponent is
\[
\Bigl[\tbinom r2-\tfrac{r(r^2-1)}3\Bigr]-(k+1)r^2+\Bigl[\tbinom r2-(r-1)r(r+1)\Bigr]+\Bigl[\tfrac r2-r(r+1)\Bigr]
=-\tfrac43r^3-(k+1)r^2-\tfrac r6=c(r,k).\qedhere
\]
\end{proof}

\begin{proposition}[Covolume in coordinates]\label{pf:sh-siegel:coordcovol}
For every word $u$ and every strict fundamental set $\mathrm{Fd}_N$ of $\Gamma_N(u)$, the set $\kappa_u(\mathrm{Fd}_N\cap
N_\infty(u))$ is measurable, of Lebesgue measure $\mathrm{cov}_N(u)/\mathrm{two}$.
\formallabel{Formalism.L2.SG.Pf.NcReal.JX.coordCovol_of}
\end{proposition}

\begin{proof}
By \lref{Formalism.L2.SG.Pf.NcReal.JX.coordCovol_of_cut}, the facts \ref{pf:sh-siegel:B1}--\ref{pf:sh-siegel:B4} imply that
the image of a strict fundamental set has Lebesgue measure $\prod_p\mathrm{cv}(p)$. The argument is the usual one for
unipotent groups: order the blocks by height. By \ref{pf:sh-siegel:B1}, left translation by $\delta\in\Gamma_N(u)$ acts in block
coordinates by $x_p\mapsto x_p+\mathrm{blk}(\delta)_p+B_p(\mathrm{blk}(\delta),x)$, a translation in the block $p$ by an amount
depending only on lower blocks; by \ref{pf:sh-siegel:B3} the translations occurring in block $p$, for fixed lower blocks, form a
coset of $\Lambda_p$. Hence, using the section of \ref{pf:sh-siegel:B2}, the image of a strict fundamental set is, block by
block and up to measure-preserving shears, a product of fundamental domains of the lattices $\Lambda_p$, and Fubini's theorem and
\ref{pf:sh-siegel:B4} give the measure $\prod_p\mathrm{cv}(p)$. \Cref{pf:sh-siegel:cvprod} concludes.
\end{proof}

\begin{proposition}[Covolume for Haar measures]\label{pf:sh-siegel:covolN}
For every word $u$, every measure $\mu$ which is Haar on $N_\infty(u)$, and every strict fundamental set $\mathrm{Fd}_N$ of
$\Gamma_N(u)$,
\[
\mu(\mathrm{Fd}_N)=\frac{\mathrm{cov}_N(u)}{\mathrm{two}}\ \mu(\mathrm{Box}(u)).
\]
\formallabel{Formalism.L2.SG.Pf.NcReal.JX.covolN_of_coordCovol}
\formallabel{Formalism.L2.SG.Pf.LsNc5.covolN_of_four}
\end{proposition}

\begin{proof}
The map $\kappa_u$ is measurable and injective on the measurable set $N_\infty(u)$ (\ref{pf:sh-siegel:B2}). Haar measures on
$N_\infty(u)$ are proportional (\lref{Formalism.L2.SG.Pf.NcReal.JX.haarUnique_holds}) and one of them has image Lebesgue
measure under $\kappa_u$ (\lref{Formalism.L2.SG.Pf.NcReal.JX.lebHaar_holds}); as $\kappa_u(\mathrm{Box}(u))=[0,1)^d$ has
Lebesgue measure $1$, the image of $\mu$ is $\mu(\mathrm{Box}(u))$ times Lebesgue measure. Since $\mu$ is concentrated on
$N_\infty(u)$ and $\kappa_u$ is injective there,
$\mu(\mathrm{Fd}_N)=\mu(\mathrm{Fd}_N\cap N_\infty(u))=\mu(\mathrm{Box}(u))\cdot\mathrm{Leb}(\kappa_u(\mathrm{Fd}_N\cap
N_\infty(u)))$, and \cref{pf:sh-siegel:coordcovol} concludes.
\end{proof}

\begin{proposition}[Reference measure of the unit box]\label{pf:sh-siegel:refbox}
For all words $u,w$, $\mathrm{ref}(u)(\mathrm{Box}(w))=T_w^{-rf}\,2^{f(rk+r^2)}$.
\formallabel{Formalism.L2.SG.Pf.NcReal.JX.refNBox_of}
\end{proposition}

\begin{proof}
At each infinite place $v$ the coordinates of $w$ are a linear image of those of $u$: $\kappa_w(s)_v=L_v(\kappa_u(s)_v)$ for a
linear map $L_v$ with $|\det L_v|=(T_w/T_u)^r$ (\lref{Formalism.L2.SG.Pf.NcReal.JX.FR.frameRow_holds}: the frames of $u$ and
$w$ differ only in the vector $z$, whose norms are $-T_u$ and $-T_w$). The image of $\mathrm{ref}_v(u)$ under the coordinates of
$u$ at $v$ is $2^{r^2+rk}T_u^{-r}$ times Lebesgue measure (\lref{Formalism.L2.SG.Pf.NcReal.JX.refN_map_kapL}). Hence the
$\mathrm{ref}_v(u)$-measure of $\{s_v:\kappa_w(s)_v\in[0,1)^{2r^2+2kr+r}\}$ is
$2^{r^2+rk}T_u^{-r}\,|\det L_v|^{-1}=2^{r^2+rk}T_w^{-r}$. Since $\mathrm{ref}(u)$ is concentrated on $N_\infty(w)$ and
$\mathrm{Box}(w)$ is $N_\infty(w)$ intersected with the product of these sets, the product over the $f$ places gives the claim.
\end{proof}

\begin{proof}[Proof of \cref{pf:sh-siegel:covolref}]
Let $w\in\{u,u'\}$ and let $\mathrm{Fd}$ be the corresponding strict fundamental set. Since $\mathrm{ref}(u)$ is Haar on
$N_\infty(w)$, \cref{pf:sh-siegel:covolN,pf:sh-siegel:refbox} give
\[
\mathrm{ref}(u)(\mathrm{Fd})=\frac{T_w^{rf}2^{-f(rk+r(r-1)/2)}\Delta^{c(r,k)}}{2^{fr(r+1)/2}}\cdot T_w^{-rf}2^{f(rk+r^2)}=
\Delta^{c(r,k)},
\]
because $\frac{r(r-1)}2+\frac{r(r+1)}2=r^2$. Hence the product of the two measures is $\Delta^{2c(r,k)}=|d_F|^{e_{DF}}$. (Only the
two fundamental-set hypotheses are used.)
\end{proof}

\subsubsection{Local volumes at the type places}

Let $v$ be a finite place of $F$, $E_v=E\otimes_FF_v$, $\OO_{E_v}$ its integral elements, $\sigma$ the conjugation of $E_v$ and
$d_v\in F_v$ an element generating the local different of $F$. For an additive character $\psi$ of $\A_F/F$ and a finite set
$\iota$, $\shsiegelsd$ is the Haar measure on $E_v^\iota$ which is self-dual for $(x,y)\mapsto\psi_v(\sum_i\Tr_{E_v/F_v}(x_iy_i))$.
Put $\alpha_a=d_v^{-(r-a)}$ for $a<r$ and $\alpha_r=1$, let $\mathcal U=\{(i,j):0\le i<j\le r\}$, and for a word $w$ let
$\mathrm{TF}_w\in M_{r+1}(E_v)$ be the Gram matrix of the frame $(e_0,\dots,e_{r-1},z_w)$ of $w$. A finite place is \emph{typed}
if it divides $g_S$; then $g=g_v\in\{g_0,g_1\}$ is the type prime below it.

\begin{definition}[Local volumes]\label{pf:sh-siegel:locvoldef}
At a typed place $v$ with type prime $g$, the \emph{local unipotent volume} is
$\mathrm{lvu}_v=(\shsiegelsd\times\shsiegelsd)(\mathcal D)$, where $\mathcal D$ is the set of pairs $(x,x')$ of elements of
$E_v^{\mathcal U}$ with $\alpha_i^{-1}\alpha_jx_{ij},\ \alpha_i^{-1}\alpha_jx'_{ij}\in\OO_{E_v}$ and
$\alpha_i^{-1}\alpha_jg^{-1}(x_{ij}-x'_{ij})\in\OO_{E_v}$; and the \emph{Gram box volume} is $\mathrm{ZZ}_v=\shsiegelsd(\mathcal
Z)$, where $\mathcal Z$ is the set of $z\in E_v^{(r+1)\times(r+1)}$ such that, with $\zeta_{ij}=g^{-1}(\alpha_i\alpha_jz_{ij}-
\mathrm{TF}_{u,ij})$, all $\zeta_{ij}\in\OO_{E_v}$ and $g^{-1}\bigl(\zeta_{ij}+\sigma(\zeta_{ji})-g^{-1}(\mathrm{TF}_{u',ij}-
\mathrm{TF}_{u,ij})\bigr)\in\OO_{E_v}$ for all $i,j$. (These are the quantities $\mathrm{lvu}_v$ and $\mathrm{ZZ}_v$ of
\cref{pf:multnc:def-bcf}.)
\formal{Formalism.L2.SG.LocVol.localVolumeUnit}\ \formal{Formalism.L2.SG.LocVol.ZZX}
\end{definition}

\begin{theorem}[Local volumes at the typed places]\label{pf:sh-siegel:voltyp}
Let $X$ be a Hida setting for a datum $D$, $u,u'$ words and $\psi$ a standard additive character of $\A_F/F$. If
$T_u\equiv T_{u'}\pmod{g_S^2}$, then for every typed place $v$ with type prime $g$,
\[
\mathrm{lvu}_v/\mathrm{ZZ}_v=g^{(r_S+1)(2r_S+3)} .
\]
\formallabel{Formalism.L2.SG.Pf.NcVal.volTypV_holds}
\formallabel{Formalism.L2.SG.Pf.NcVal.volTypV_of_box}
\formallabel{Formalism.L2.SG.Pf.NcVal.volTypV_of_herm}
\formallabel{Formalism.L2.SG.Pf.NcVal.volTypV_of_cnt}
\formallabel{Formalism.L2.SG.Pf.NcVal.volTypV_of_steps'}
\formallabel{Formalism.L2.SG.Pf.NcVal.volU_steps_of}
\end{theorem}

\begin{proof}
Write $r=r_S$, $O=\OO_{E_v}$, $|\cdot|_v$ for the normalised absolute value of $F_v$.

\emph{Step 0: arithmetic of the typed places.} The type prime $g$ splits completely in $E$ (the number of primes of $E$ above
$g$ is $[E:\Q]$, \lref{Formalism.L2.SG.Pf.NcVal.VolUT.gAt_split}), so the place $v$ of $F$ above $g$ has ramification index
and residue degree $1$ and splits in $E$: $E_v\cong F_v\times F_v$, with $\sigma$ exchanging the factors. Its residue field has
$g$ elements and $g\notin v^2$, so $|g|_v=g^{-1}$; and $v$ does not divide the different of $F$, so $d_v$ is a unit, $|d_v|_v=1$,
and all $\alpha_a$ are units (\lref{Formalism.L2.SG.Pf.NcVal.VolUT.typ_arith}).

\emph{Step 1: the self-dual measure.} Since $\psi$ is standard, its local component at $v$ has conductor the inverse different
(\lref{Formalism.L2.SG.Pf.conductor_inverse_different}); as $d_v$ is a unit, the dual of $O$ for $(x,y)\mapsto
\psi_v(\Tr(xy))$ is $d_v^{-1}O=O$ (\lref{Formalism.L2.SG.Pf.NcVal.VolUTr.annX_of}). Hence $O^\iota$ is self-dual, the
self-dual measure exists and is a Haar measure, and $\shsiegelsd(O^\iota)=1$ for every finite $\iota$
(\lref{Formalism.L2.SG.Pf.NcVal.VolUL.sd_facts}).

\emph{Step 2: scaling.} For $b\in E_v^\iota$ and measurable $S$,
$\shsiegelsd\{w:\ g^{-1}(w-b)\in S\}=g^{-2|\iota|}\shsiegelsd(S)$: the map $w\mapsto g^{-1}(w-b)$ is affine, with linear part
multiplication by $g^{-1}$ on $E_v^\iota\cong F_v^{2|\iota|}$, of module $|g|_v^{-2|\iota|}=g^{2|\iota|}$
(\lref{Formalism.L2.SG.Pf.NcVal.VolUL.volTAff_of}). Multiplying the coordinates by the units $\alpha_i\alpha_j$ preserves
$\shsiegelsd$ (\lref{Formalism.L2.SG.Pf.NcVal.VolUL.volTLocPres_of}).

\emph{Step 3: hermitian count.} Let $c\in M_{r+1}(O)$ with $c_{ji}=\sigma(c_{ij})$, and
$\mathcal W(c)=\{Y\in O^{(r+1)\times(r+1)}:\ g^{-1}(Y_{ij}+\sigma(Y_{ji})-c_{ij})\in O\ \text{for all }i,j\}$. We claim
$\shsiegelsd(\mathcal W(c))=g^{-(r+1)^2}$. As $c$ is hermitian, the condition for $(j,i)$ is the conjugate of the condition for
$(i,j)$. For $i<j$, given $Y_{ij}\in O$ the condition says $Y_{ji}\in\sigma(c_{ij}-Y_{ij})+gO$, a set of measure $g^{-2}$. For
$i=j$ the condition says $\Tr(Y_{ii})\in c_{ii}+gO_v$; in the decomposition $O=O_v\times O_v$ (Step 0) the trace is
$(a,b)\mapsto a+b$ and $\shsiegelsd$ is the product of the Haar measures of $O_v$ of total mass $1$, so this set has measure
$|g|_v=g^{-1}$. By Fubini,
$\shsiegelsd(\mathcal W(c))=(g^{-2})^{r(r+1)/2}(g^{-1})^{r+1}=g^{-(r+1)^2}$. (The formal proof writes $\mathcal W(c)$ as a
translate of the image of $O^{(r+1)\times(r+1)}$ under an explicit linear map built from an element $\theta\in O$ of trace $1$ and
the element $\theta-\sigma(\theta)$, of determinant module $g^{-(r+1)^2}$:
\lref{Formalism.L2.SG.Pf.NcVal.volUHermBox_holds} and \lref{Formalism.L2.SG.Pf.NcVal.VolUHm.volTHermCnt_of_box}.)

\emph{Step 4: the Gram box.} The frames of $u$ and $u'$ differ only in the last vector, so $\mathrm{TF}_{u'}-\mathrm{TF}_u$ has the
single non-zero entry $T_u-T_{u'}$ in position $(r,r)$; as $g^2\mid g_S^2\mid T_u-T_{u'}$, $c=g^{-1}(\mathrm{TF}_{u'}-\mathrm{TF}_u)$
is hermitian with integral entries (\lref{Formalism.L2.SG.Pf.NcVal.volTTFCong_of}). By definition
$z\in\mathcal Z$ if and only if $g^{-1}((\alpha_i\alpha_jz_{ij})_{ij}-\mathrm{TF}_u)\in\mathcal W(c)$. By Steps 2 and 3,
\[
\mathrm{ZZ}_v=g^{-2(r+1)^2}\,\shsiegelsd(\mathcal W(c))=g^{-3(r+1)^2}.
\]
\emph{Step 5: the unipotent volume.} As the $\alpha_a$ are units, $\mathcal D=\{(x,x'):\ x\in O^{\mathcal U},\ g^{-1}(x'-x)\in
O^{\mathcal U}\}$ (\lref{Formalism.L2.SG.Pf.NcVal.volTUOEq_of}), and by Fubini and Steps 1--2,
$\mathrm{lvu}_v=\shsiegelsd(O^{\mathcal U})\,g^{-2|\mathcal U|}=g^{-r(r+1)}$
(\lref{Formalism.L2.SG.Pf.NcVal.VolUL.volTProdU_of'} and \lref{Formalism.L2.SG.Pf.NcVal.volU_two_card_uidx}).

\emph{Step 6.} $\mathrm{lvu}_v/\mathrm{ZZ}_v=g^{3(r+1)^2-r(r+1)}=g^{(r+1)(2r+3)}$.
\end{proof}


%==============================================================================
\subsection{The length formula on the second branch}\label{sec:pf:sh-siegel:branch}

Let $D$ be a datum and $u$ a word, with member forms at level $0$: dominant weights $\lambda_0(\tau)$ ($\tau$ an embedding type),
the space of forms of $u$ at level $0$, and for every place $w$ of the coefficient field $\mathit{LL}$ the localised lattice
$S_0(w)$ of forms and the lattices $B_0(w,v)$ (for the indices $v$ of the level). For $\tau=\tau(w,v)$ put
$b(\tau)=-\mathrm{ds}(\lambda_0(\tau))$ (minus the determinant shift of the weight) and $|\lambda_0(\tau)|$ the sum of its entries.
For $\OO_w$-lattices $A,B$ let $\len_w(A,B)$ be the length of the $\OO_w$-module $A/(A\cap B)$, and for $n\in\Z$ let
$T_u^n\cdot S_0(w)$ be the image of $S_0(w)$ under multiplication by $T_u^n$. A place $w$ \emph{lies over} $T_u$ if its residue
characteristic is $T_u$; it is \emph{seen} by $u$ if the forms of $u$ record it; $e_w$ is its ramification index; and $d^S_u$ is
the size of $u$.

\begin{definition}[The second-branch length]\label{pf:sh-siegel:branchdef}
The \emph{second-branch inclusion length} $\mathrm{SB}(D,u)$ is: for every place $w$ over $T_u$ which is seen by $u$ and every
$v$ with $b=b(\tau(w,v))\ge0$, $T_u^b\cdot S_0(w)\subseteq B_0(w,v)$, and $\len_w(B_0(w,v),T_u^b\cdot S_0(w))$ is a natural number
$n$ with
\[
n=e_w\cdot\frac{d^S_u\bigl(|\lambda_0(\tau)|+k\,b\bigr)}{k},\qquad\tau=\tau(w,v).
\]
The \emph{level-zero hypotheses} on $(D,u)$ are: the numerical conditions $2fN_m+1<T_u$, $4a_{\max}(u)<T_u$, the alcove condition
and $b(\tau)\ge0$ for the top types $\tau$; and the concrete forms at level zero (compatibility of the forms with the twists
along the family, and the admissibility of the size data $\mathrm{adm}_S(D,u)$).
\formal{Formalism.L2.GlobalTheoryTheory.LocalLatticeBases.SecondBranchInclusionLength}\
\formal{Formalism.MainNotions.RT.SecondBranchIntegralityUnderForms}
\end{definition}

\begin{theorem}[The second-branch length formula]\label{pf:sh-siegel:secondbranch}
For every datum $D$ and every word $u$ satisfying the level-zero hypotheses, $\mathrm{SB}(D,u)$ holds.
\formallabel{Formalism.L2.GlobalTheoryTheory.LemB.second_branch_all_members}
\formallabel{Formalism.L2.GlobalTheoryTheory.LemB.second_branch_of_rationality}
\formallabel{Formalism.L2.GlobalTheoryTheory.LemB.second_branch_of_surjectivity}
\formallabel{Formalism.L2.GlobalTheoryTheory.LemB.second_branch_of_coordinates_and_count}
\end{theorem}

The proof is a computation of lengths in coordinates. Let $\Delta(\lambda)\subseteq\C[x_{ab}]$ ($0\le a,b<N_0$) be the span of the
products of leading principal minors attached to a dominant weight $\lambda$ (the highest weight module of $\lambda$), and
$\Delta_\Z(\lambda)$ its integral form; the forms of $u$ at level $0$ are realised in the tensor product
$\bigotimes_\tau\Delta(\lambda_0(\tau))$. Let $\mathrm{Rat}(\lambda_0)$ be the span of the forms with rational coefficients
(over the coefficient field) and $\mathrm{CSpan}(\lambda_0)$ the complex span of the pure tensors of elements of the
$\Delta(\lambda_0(\tau))$.

\begin{lemma}[Rationality]\label{pf:sh-siegel:rat}
\begin{enumerate}[label=(\alph*)]
\item $\mathrm{Rat}(\lambda_0)\subseteq\mathrm{CSpan}(\lambda_0)$.
\item For every family $\gamma=(\gamma_\tau)$ in $\GL_{N_0}(E)$, with images $o_\tau$ at the infinite places, the substitution
$P\mapsto\prod_\tau\det(o_\tau)^{s(\tau)}\cdot P(o\,\cdot)$, $s(\tau)$ the determinant shift of $\lambda_0(\tau)$, maps
$\mathrm{Rat}(\lambda_0)$ into itself.
\end{enumerate}
\formallabel{Formalism.L2.GlobalTheoryTheory.LemB.ratInC_holds}
\formallabel{Formalism.L2.GlobalTheoryTheory.LemB.ratGL_holds}
\end{lemma}

\begin{proof}
(a) $\mathrm{Rat}(\lambda_0)$ is the additive group generated by the multiples of pure tensors $\bigotimes_\tau f_\tau$ of
Schur--Weyl images of integral elements; each $f_\tau$ lies in $\Delta(\lambda_0(\tau))$
(\lref{Formalism.L2.G.Model.sw_weyl_map}), so each generator lies in $\mathrm{CSpan}(\lambda_0)$, and the claim follows by
induction on the generation (closure under sums and negatives). (b) The images $o_\tau$ are the images of matrices $M_\tau$ over
the coefficient field with non-zero determinant. On a generator, the substitution acts as a scalar times the tensor product of
the row actions of the $M_\tau$ on the factors (\lref{Formalism.L2.GlobalTheoryTheory.LocalLatticeBases.substitution_tensor}); the row
action of an invertible matrix preserves the rational span of each $\Delta(\lambda_0(\tau))$
(\lref{Formalism.L2.GlobalTheoryTheory.LemB.rat_rowAct_mem}), and tensor products of rational elements are rational
(\lref{Formalism.L2.GlobalTheoryTheory.LemB.rat_ten}). Induction on the generation concludes.
\end{proof}

\begin{proof}[Proof of \cref{pf:sh-siegel:secondbranch}]
\emph{Step 1: surjectivity of coordinates.} For a place $w$ over $T_u$ seen by $u$, an index $v$ with $b\ge0$, and families $B,B'$
of local bases (indexed by the classes at level $0$) of the lattices of $u$ at level $0$ and of the twisted lattices at $v$, the
\emph{surjectivity condition} $\mathrm{Onto}(w,v,B,B')$ says that the images of the integral forms under the coordinate maps
attached to $B$ contain, respectively, the $\det^{s}$-twisted substitution of the integral forms by $B^{-1}B'$ (inside the
coordinates of $B_0(w,v)$) and the $T_u^b$-multiples of the integral forms (inside the coordinates of $T_u^b\cdot S_0(w)$). By
\cref{pf:sh-siegel:rat} and \lref{Formalism.L2.GlobalTheoryTheory.LemB.ontoAll_of}, $\mathrm{Onto}(w,v,B,B')$ holds for all
such data.

\emph{Step 2: coordinates with determinant $T_u$.} Given the surjectivity of Step 1, there are local bases $B,B'$ such that
$\det(B_c^{-1}B'_c)$, at the embedding $\tau(w,v)$, has the same $w$-adic valuation as $T_u$ for every class $c$
(\lref{Formalism.L2.GlobalTheoryTheory.LemB.second_branch_coordinates_of_surjectivity}, using the vanishing of the valuation of
basis changes at level zero, \lref{Formalism.L2.GlobalTheoryTheory.LemB.V0.det_basis_change_unit_zero}, and the second
unramified condition of the host, \lref{Formalism.L2.Host.second_unramified_condition_holds}).

\emph{Step 3: the count.} Under the forms at level zero and $\mathrm{adm}_S(D,u)$, $d^S_u$ equals the number of classes at level $0$
times $\prod_\tau\rank_\Z\Delta_\Z(\lambda_0(\tau))$ (\lref{Formalism.L2.GlobalTheoryTheory.LemB.duCount_holds}).

\emph{Step 4: lengths.} With the coordinates of Step 2, an injective semilinear coordinate map identifies $B_0(w,v)$ and
$T_u^b\cdot S_0(w)$ with products over the classes $c$ of, respectively, the $\det^s$-twisted substitution by $B_c^{-1}B'_c$ of the
integral forms and the $T_u^b$-multiples of the integral forms (injectivity of the coordinates,
\lref{Formalism.L2.GlobalTheoryTheory.LemB.crdL_inj}, the two inclusions of Step 1 and the reverse inclusions
\lref{Formalism.L2.GlobalTheoryTheory.LemB.crdL_Bw_le} and \lref{Formalism.L2.GlobalTheoryTheory.LemB.crdL_tL_le}). Lengths are
preserved by injective semilinear maps and add over products
(\lref{Formalism.L2.GlobalTheoryTheory.LemB.len_quot_semilinear} and \lref{Formalism.L2.GlobalTheoryTheory.LemB.len_pi_sum}). For
each class the length $n_c$ of the quotient satisfies
$k\,n_c=e_w\prod_\tau\rank_\Z\Delta_\Z(\lambda_0(\tau))\cdot(|\lambda_0(\tau(w,v))|+kb)$
(\lref{Formalism.L2.GlobalTheoryTheory.LemB.cls_len_D}, computed in a homogeneous basis of the modules $\Delta_\Z(\lambda_0(\tau))$
from the valuation of $\det(B_c^{-1}B'_c)$ of Step 2). Summing over the classes and using Step 3 gives the length
$e_wd^S_u(|\lambda_0(\tau)|+kb)/k$, which is the seen-place length formula
(\lref{Formalism.L2.GlobalTheoryTheory.LemB.lenLen2_of_coord} and \lref{Formalism.L2.GlobalTheoryTheory.LemB.lenLen2_of_coordOnto}).

\emph{Step 5: inclusion.} The seen-place length formula, together with the inclusion $T_u^b\cdot S_0(w)\subseteq B_0(w,v)$ at level
zero (\lref{Formalism.L2.GlobalTheoryTheory.LocalLatticeBases.second_branch_incl_length_level_zero'}), gives $\mathrm{SB}(D,u)$
(\lref{Formalism.MainNotions.lemmaB_second_branch_of_length_formula}).
\end{proof}

